\section{Weighted comparison budgets}\label{sec:budgets}

The purpose of this section is to turn optimized pressure comparisons into
uniform second-moment budgets for the genuine paths. The global
comparison gives both parts of \eqref{eq:C2} directly. To prove the
remaining low-mass part of \eqref{eq:C1}, we first obtain coarse
comparability for auxiliary minimizers and then strengthen it to the
accuracy required by a local pressure comparison. The local auxiliary
optimizers are used only to prove those budgets.  Their density constraints
will change during the argument, so their regularity and concentration
estimates are established anew before any one-site closure is applied.
All constants and all finite iteration counts are fixed before the size
limit.  Parameter choices are recorded where they first become necessary;
the final order is collected in Section~\ref{sec:parameters}.

\subsection{Averaging targets and the global comparison}

Unlettered equation numbers in references to source estimates are those
of \cite{PE}. Its numbered statements use the source numbering.
The estimates of this section carry labels (C1)--(C20).
Whenever the density constraints differ from those in \cite{PE}, we
state the replacement analytic inputs and repeat the relevant proof
argument below.

Write \(A_v=B_v-2q(v)S_v+q(v)^2\) for the actual untilted second moment about the reference. Fix \(k_c=0.03\). In what follows take \(s_*=1/2-\tau\), \(s_0=s_*-\tau_0\), with
\[
 0<\tau_0\ll\eta\ll\tau\ll c_G,\qquad \tau\ll 10^{-4}
\]
(\(1-2s_0\) can in particular be made arbitrarily small relative to \(c_G/12\)). The final-strength exponent \(p_*\) can be taken sufficiently small after these choices.


\begin{proposition}[Uniform weighted comparison budgets]\label{bud:targets}

The averaging bounds needed are, on all segments of Proposition~\ref{prop:increment}, uniformly,
\[
 \mathbf E_\gamma\left[\max_{U\lesssim u\lesssim1}
       u\int_{v\asymp u} A_v\,dv\right]\lesssim_\circ U^4,              \tag{C1}\label{eq:C1}
\]
where one can use ordinary dyads with fixed comparable enlargements, and on size paths in addition
\[
 \int (v^2+U^2)\mathbf E_\gamma A_v\,dv \le x^{4+c},\qquad
 f_{\rm spin}(K)+\int (K-h)^2/(4b)- f_m^0 \le x^{4+c}.                  \tag{C2}\label{eq:C2}
\]
The second assertion of \eqref{eq:C2} is pointwise, with some \(c>0\) independent of sufficiently small \(p_*\).

\end{proposition}\begin{proof}[Reduction to localized comparisons]


We first obtain the size-path bounds from the global comparison. Fix a
direct endpoint and an unseeded, unmasked path whose global parameter
is smaller by a fixed \(\delta>0\). To reach the endpoint, remove
\(\delta c_g\) from the matrix clock and add \(\delta c_gq+d\) to
the field clock. Both paths are admissible. Completing the square
along their affine interpolation and using \eqref{eq:G1} bounds
\[
 R=f_m^0-\int bq^2/4+\int_{v\ge x}d q/2
\]
below by \(f_{\rm spin}(K)-O_\delta({\cal B})\); the corresponding upper bound follows by the trial replacing the remaining matrix to reach \(K\). Here, with subpower losses harmless,
\[
 {\cal B}\lesssim_\circ x^{4+c_G}+
       \lambda^{1+2\eta}U^5+\chi_e^2\lambda^{1+2\eta}U^4+\chi_{\rm lin}^2 x^6/\lambda.
\]
The lower comparison costs \(\int d^2/(\delta c_g)\), and the upper
comparison costs \(\int d^2/b\). Omitting the mask costs
\(O(\lambda U^2x^3)\). The linear restoration is independent of
\(\gamma\), so \eqref{eq:P10} can be integrated between the global
parameter endpoints. For some fixed \(b_s>0\), this gives
\[
 \int(v^2+U^2)\,\mathbf E_{\gamma_0} A_v\,dv
       \lesssim x^4\lambda^{-2s_0+b_s}
\]
also when the global interval of averaging is any fixed interior interval leaving crossing room. We can choose \(b_s,\tau_0\) sufficiently small relative to \(\eta\), with further room to absorb errors \(x^4u\lambda^{k_c/3-1-o(1)}\) on the right. At final strength with edge off this argument and \eqref{eq:P11} give both parts of \eqref{eq:C2}, with savings independent of further decreases in \(p_*\): even the differentiated estimate uses only \(O_\circ(x^{4+c_G}/\lambda+\lambda^{2\eta}U^5+x^6/\lambda^2)\) plus the switch cost from \eqref{eq:P11} divided by \(\lambda\).

For \eqref{eq:C1}, the remaining task is a direct-path dyad \(u\) of
\eqref{eq:P3} with \(U\lesssim u\le\lambda^{b_\ell}\), where
\(b_\ell>0\) is sufficiently small relative to \(b_s\). Bounded
overlaps already suffice when \(U\) is large. Fix all parameters
other than \(\gamma_0,\gamma_u\), and start at \(\gamma_u=0\).
At this endpoint we minimize over an auxiliary quantile \(r\): the
matrix removal is \(C(r)>0\), the field addition is \(C(r)r\),
and the penalty is \(\int C(r)r^2/4\). We will specify its density
constraints below. Writing \(X=x/u\), it is enough to prove the
following two assertions for the final constraints:

\begin{enumerate}[label=(\roman*)]

\item averaged excess of this optimized value over \(f_m^0\) at most \(O(x^4 u\lambda^{1-2s_0})\);
\item for the used optimizers, on \(\operatorname{supp}(c_u)\), \(c_u=-\partial_{\gamma_u}b\), one has \(C(r)\gg c_u\) and
\[
 \mathbf E_{\gamma_0}\left[\sup |r-q|^2\right]
           \lesssim (u X^2\lambda^{-s_0})^2 .                            \tag{C3}\label{eq:C3}
\]

\end{enumerate}


To see this, compare a minimizing auxiliary path to the actual endpoint
\(\gamma_u=1\). Starting at that actual path, the square interpolation
removes \(C(r)-c_u\) and adds \(C(r)r-c_uq\). The increase of the
restored pressure \(R\) between the two endpoints is therefore at most
the optimizer excess plus
\(\int C(r)c_u(r-q)^2/[4(C(r)-c_u)]\).
Averaging this bound controls the nonnegative derivative in
\eqref{eq:P10}; summing over dyads costs at most a logarithm.
Before differentiating on a smoothed size path, transfer the same
pressure difference by \eqref{eq:P11}. Its cost is \(O(x^{5+c})\),
since only the restoration \(-\int bq^2/4\) affects the difference.
Fixed rescalings of mass change these bounds only by bounded factors.

The second requirement is an estimate for every optimizer outside one
common null set of the global parameter.  It is enough to construct a
measurable scalar majorant for all these optimizers; no choice of an
optimizer needs to be integrated.  The remaining subsections establish
these two local assertions and hence complete \eqref{eq:C1}. Both
parts of \eqref{eq:C2}, with a saving independent of the local
strengthening depth, have already been proved above.
\end{proof}



\subsection{Local auxiliary comparisons and coarse profiles}

For the local argument normalize each direct path to length 1, reusing notation for its matrix and field coefficients multiplied by \(\ell^2\), the scaled scalar clock and moments, and writing \(v\) for normalized mass. Thus \(q(v)\) here is the old \(q(\ell v)\), and \(q\) and \(q^{-1}\) have bounded derivatives and positive slopes on fixed initial intervals.
All pressure bounds (with their corrections) change by bounded normalization factors. The scale \(u\) of \eqref{eq:C3} need not be rescaled: its relevant mass range is still contained in a fixed compact positive band of scale \(u\). Fix such a dyad, with \(\gamma_u=0\). In particular
\[
 x\lambda^{-s_*/2}\lesssim u\le\lambda^{b_\ell},\quad X=x/u,\quad
 \Lambda=\lambda^{k_c},\quad A=X^2\Lambda^{-1/3},\quad a_{\rm goal}=X^2\lambda^{-s_0}.
                                                                    \tag{C4}\label{eq:C4}
\]
Thus \(U=x\lambda^{-s_*/2}\) for large sizes here. Use an arbitrarily small fixed \(\zeta>0\) (depending on previously required power margins including \(p_*,b_\ell\), decreased as needed). Apply \cite[Proposition~6.6]{PE} with
\[
 d_{\rm PE}=c_d u\Lambda^{1/3},\qquad E=x^{4-\zeta}d_{\rm PE},\qquad
 C(r)=C_0 f(C_0 r/d_{\rm PE}^3),\qquad C_0=u^2\Lambda ,
\]
i.e. at strength \(C_0\); \(f\) here is the PE taper of (6), constant 1 near 0. Choose the fixed \(c_d\) sufficiently large to have this constant profile on \([0,c_I u]\) for an arbitrarily large fixed \(c_I\), in particular well beyond the compact range of scale \(u\) where \eqref{eq:C3} will be tested.

PE power slack (its small parameter in (3), and \(\delta\)) is chosen sufficiently small relative to \(\zeta\). Write \(k(r)\) for the PE floor (5), \(v_*=(mE)^{-1}\asymp x^\zeta u A\). These parameters satisfy (3) since \(X\lesssim\lambda^{s_*/2}\), and \(C_0\le d_{\rm PE}^2=o(1)\). Conditions (4)--(6) with base \(T_m\to T>1\) hold by \eqref{eq:P3}--\eqref{eq:P9}; all masses are changed, with no radius wall, and a fixed positive covariance margin throughout. Use the common-base version explained at \eqref{eq:G1}. The optimized excess over the actual untilted pressure is thus \(\lesssim x^{4-2\zeta}u\Lambda^{1/3}\).

The modified optimizations below, at this same strength and with the same objective \(f_m(b-C(r),h+C(r)r)+\int C(r)r^2/4\), use inverse constraints on \([0,1]\)
\[
 \bar k(r)\le\alpha'(r)\le M_j,\qquad
 \bar k(r)=k(r)+\kappa\,{\bf1}_{r\in I},\qquad
 0\le\kappa\lesssim\lambda^{2s_0}\Lambda^{-2/3}.                        \tag{C5}\label{eq:C5}
\]
Here \(I/u\) is a fixed closed subinterval of \((0,c_I)\), and
\(M_j=2^{j+1}x^{-100}\), with only finitely many fixed stages
\(j\ge-1\). The inverse distribution \(\alpha\) maps onto \([0,1]\);
its inverse \(r=\alpha^{-1}\) is the forward quantile. The initial
choice \(\kappa=0,\ j=-1\) contains the original PE class.
All these constraints are deterministic and independent of
\(\gamma_0\). Compactness and the segment continuity in
\eqref{eq:A1} give a minimizer. In the local argument, the symbols
\(S,B,D,X_v,W_v\) refer to the changed path of any such minimizer,
unless stated otherwise.

\begin{lemma}[Modified optimizer geometry]\label{bud:geometry}
For every minimizer in \eqref{eq:C5}, the first-variation potential can be normalized so that its derivative is
\[
 {\cal H}'=C(r)\big[(2-p)r-2(1-p)S-p B/r\big],\qquad p=-r C'/C .
                                                                    \tag{C6}\label{eq:C6}
\]
The minimizing density equals the cap $M_j$ where $\mathcal H<0$
and the floor $\bar k$ where $\mathcal H>0$; the remaining points are
contacts. Each cap has radius width at most $M_j^{-1}$. Outside the
floor, apart from a null set of contact exceptions,
$|S-r|\lesssim_\circ D+O(x^{90})$. At contacts where $C$ is constant, $S=r$; on every cap contained
in that constant region, $|S-r|\le M_j^{-1}$.
These conclusions are the parts of the geometry of
\cite[Lemma~3.3, equations~(9)--(10)]{PE} that will be used below.

\end{lemma}\begin{proof}

The first variation minimizes a linear functional of \(\alpha'\),
subject to its prescribed lower and upper bounds and total mass one.
The multiplier for total mass adds a constant to the primitive
\({\cal H}\). With this choice, the minimizing density equals
\(M_j\) where \({\cal H}<0\), and \(\bar k\) where
\({\cal H}>0\). These are the cap and floor regions, respectively;
the remaining points are contacts. Differentiation of inverse monotone
maps is justified as at \eqref{eq:G1}; the radii at which a jump is
crossed form a null set.

The monotonicities of \(S,B,p\), together with the zero initial wall,
give the endpoint signs, almost-everywhere contact relations, and the
cap bounds in PE (9)--(10). Each cap has radius width at most
\(M_j^{-1}\). We call cap and contact points \emph{nonfloor points},
discarding the null set of contact exceptions. There
\(|S-r|\lesssim_\circ D+O(x^{90})\); on a component in the constant
part of \(C\), the stronger bound is \(|S-r|\le M_j^{-1}\).
Also \(S\le r\) at almost every contact. On a negative cap, \(S\)
cannot exceed its right endpoint radius: otherwise \(B\ge S^2\)
and monotonicity make \({\cal H}'<0\) on a terminal portion of the
cap, contradicting its endpoint sign. At the outer endpoint the same
argument uses \(S\le1\).

\end{proof}


Put
\[
 y_c=x^{-10\zeta}\big(X\Lambda^{-1/6}+X^{3/4}\lambda^{0.12}\big),\qquad w_c=y_c u .
\]
Here and in this local argument \(\delta\) on a PE power denotes \(\delta_{\rm PE}\) as recalled at \eqref{eq:G1}. In particular \(a_{\rm PE}\asymp x^{1-\zeta}\); \(w_c/u\ll\sqrt\Lambda\), \(x/w_c\) is power small, and
\(w_c\ge m^\delta\sqrt m E/\sqrt{C_0}\). Indeed \(X\lesssim\lambda^{s_*/2},\ \lambda\le x^{p_*}\) allow the prefactor \(x^{-10\zeta}\) in the first assertion if \(\zeta\) is small enough. Also \(C(w_c)=C_0,\ C_0 w_c^2\gtrsim x^{-18\zeta}mE^2\).
\Needspace{4\baselineskip}
\begin{proposition}[Comparability and the calibrated initial slope]\label{bud:coarse-profile}
For every optimizer in \eqref{eq:C5}, uniformly,
\[
 r(v)\asymp v\ (w_c\le v\le1),\qquad
 \sup_{c u\le v\le C u}|r(v)-q(v)|=o(u)                                 \tag{C7}\label{eq:C7}
\]
for each fixed \(0<c<C\) (fixed numeric multiples, not the removal profile). \end{proposition}

The proof has two uses for concentration estimates. During the
comparability contradiction they are available only outside a possible
irregular region; after comparability is proved they apply with the
threshold $w_c$. We first give this conditional argument, without assuming
the conclusion of \eqref{eq:C7}.

\subsubsection{Analytic estimates conditional on outer regularity}

Fix an optimizer of \eqref{eq:C5} and a scale $w\ge w_c$ with
$w=o(u)$. Suppose that, for fixed positive constants, its quantile is
comparable to mass above $C_{\rm reg}w$. Thus $r(v)\asymp v$ there;
by monotonicity $r(v)=O(v+w)$ everywhere. We call bands lying beyond a
sufficiently large fixed multiple of this threshold regular. The scale
bounds used in the argument are
\[
 w\ge m^\delta\sqrt m E/\sqrt{C_0},\qquad
 C(w)w^2\gtrsim m^{1+\delta}E^2,\qquad m^{2\delta}v_*\ll w.
\]
In both applications below these follow from $w\ge w_c$. The assumption
$w=o(u)$ places the enhanced-floor patch $I$ in the regular region.
All assertions in the next lemma are conditional on this fixed optimizer
and these outer bounds. No minimizer is selected as the parameters vary.

\begin{lemma}[Moments under outer regularity]\label{bud:outer-moments}
For every fixed $2\le p<\infty$, at every specified revelation cut and fork,
\[
 \|X_v\|_p\lesssim_p\sqrt{v+w},\qquad \|Q_v\|_p\lesssim_p v+w              \tag{C8}\label{eq:C8}
\]
where $Q_v$ is the overlap of a pair splitting at mass $v$.
On a regular small radius band $J\asymp v$,
\[
 \left\|\int_J W_{\alpha(r)}\,dr\right\|_{p/2}
 \lesssim_{p,\circ}x^6/C(v).
\]
When a preceding regular projection interval of comparable scale fits,
the synchronized two-branch residual has norm
$O_{p,\circ}(x^3/(v\sqrt{C(v)}))$.
The constants are uniform over optimizers with the stated outer bounds.
\end{lemma}
\begin{proof}
\leavevmode\par
\paragraph{Raw moments under outer regularity.} At regular dyads of scale \(v\) (mass and radius comparable beyond a sufficiently large fixed multiple of \(w\)) one has field slope at least \(C(v)/O(\log(1/x))\) in radius, and
\[
 m C(v)k(v)v^4\gtrsim m^\delta C(v)v^2/(mE^2)
\]
diverges polynomially. We have \(S=O(v+w)\) by looking ahead to a nonfloor point (a band of mass comparable to regular \(v\) has floor mass \(o(v)\)) or by monotonicity below regular scales. We will prove the raw bounds stated above by an absorption argument, following PE (12).

Fix the optimizer, all path parameters, and \(l\ge1\). Define
\[
 Q_{\rm raw}=\max\left\{1,
   \sup_t\frac{\|X_t\|_{4l}}{\sqrt{v(t)+w}}\right\},
\]
where \(t\) ranges over deterministic revelation cuts, including the
specified intermediate cuts, and \(v(t)\) is their mass. At fixed size
this is a finite deterministic number: \(w>0\) and \(|X_t|\le1\).
We will bound it uniformly, without assuming an upper bound on the
mass of a cap. On a regular small radius band \(J\asymp v\), we claim
\[
 \left\|\int_J W_{\alpha(r)}\,dr\right\|_l
    \lesssim_{l,\circ} Q_{\rm raw}^2 x^6/C(v).                            \tag{C9}\label{eq:C9}
\]
Fork an independent continuation at each radius \(r\). Decompose both
residuals into future martingale increments, using the weighted Hessian
representation in Lemma~\ref{global:residual-blocks}. Retaining
its mass weight is necessary here. First use dyadic distances from a
sufficiently small inverse power of size up to a fixed fraction of
\(v\). Then use doubling scales up to a fixed positive radius and the
terminal layer. For a pair of layers with largest scale \(z\gtrsim v\),
the Hessian weight is \(O(z)\); for two short layers it is \(O(v)\).

For short layers, let \(h'\) be the starting distance of the later one.
Unless both are the first increment, \eqref{eq:A3} bounds the squared
product by the earlier squared anchor times
\(x^{6-\epsilon}/(v^2C(v)h')\), apart from a negligible tail.
Here one first uses arbitrarily high fixed moments for the normalized
anchor, then truncates; boundedness of the actual increments controls
the discarded tail. The shifted-grid and dual conditional-expectation
estimates used for \eqref{eq:G4} bound the integral of the earlier
squared increment by \(O_l(h'Q_{\rm raw}^2v)\) in \(L^l\).
Thus the three factors in a short-layer contribution are
\[
 \underbrace{v}_{\text{Hessian weight}}\,
 \underbrace{\frac{x^{6-\epsilon}}{v^2C(v)h'}}_{\text{projection}}\,
 \underbrace{h'Q_{\rm raw}^2v}_{\text{integrated anchor}}
 =\frac{x^{6-\epsilon}Q_{\rm raw}^2}{C(v)}.
\]
Two first increments are negligible by the square-sum estimate and the
choice of their inverse-power length.

For a larger layer of scale \(z\), there is a preceding regular
interval of length and mass comparable to \(z\). Projection costs
\(x^{6-\epsilon}/(z^3C(z))\), and the squared opposite anchor has
\(L^l\) norm \(O_l(Q_{\rm raw}^2z)\). Its integrated contribution is
therefore at most
\[
 z\,\frac{x^{6-\epsilon}}{z^3C(z)}\,
 Q_{\rm raw}^2z\,|J|
 =\frac{x^{6-\epsilon}Q_{\rm raw}^2|J|}{zC(z)}
 \lesssim\frac{x^{6-\epsilon}Q_{\rm raw}^2}{C(v)}.
\]
The last step uses \(|J|=O(v)\) and \(zC(z)\gtrsim vC(v)\),
also at the fixed outer cutoff. Summing the logarithmically many layer
pairs proves \eqref{eq:C9}.

By \eqref{eq:A1} and the floor lower bound, the centered average of
\(L\) over a future radius window of length comparable to \(v\)
has \(L^{2l}\) norm at most \(\varepsilon_m vQ_{\rm raw}\), where
\(\varepsilon_m\to0\) with a power saving. At scale \(v\), its bound
is \(O_\circ([mC(v)k(v)v^4]^{-1/2})\); the polynomial divergence
proved above absorbs the subpower loss. Conditional expectation
back to an earlier cut bounds \(|X|^2/2\) by this average. Its
deterministic mean is \(O(v)\). For a cut inside the nonregular region,
choose the window at a fixed larger regular scale comparable to \(w\);
for a regular small cut, choose it at a fixed multiple of that cut's
scale. Macroscopic cuts use \(|X|\le1\). Taking the supremum over
these deterministic cuts gives
\[
 Q_{\rm raw}^2\le C_l+\varepsilon_m Q_{\rm raw}.
\]
Finiteness of \(Q_{\rm raw}\) now allows absorption and proves its
uniform bound. For a small fork, project both branches into the same
type of later window and telescope the remaining synchronized layers,
as in PE (16). If \(v\) is the window scale, the residual is
\(O_{p,\circ}(x^3/(v\sqrt{C(v)}))\ll v\). This proves
\eqref{eq:C8}, and the PE (16) postprojection bounds whenever the
required preceding regular intervals fit.

\end{proof}

\Needspace{4\baselineskip}
\begin{lemma}[Second moments with an enhanced floor]\label{bud:outer-second}
Under the same outer-regularity hypotheses, on regular mass dyads, with
$B_f=X^{3/2}\lambda^{s_0/2}\Lambda^{-5/12}$,
\[
 (\operatorname{avg}D^2)^{1/2}
 \lesssim_\circ (x^6/(C(v)v))^{1/3}+x^3/(v\sqrt{C(v)})
            +m^{\delta/2}v_*+{\bf1}_{v\asymp u}\,u B_f .               \tag{C10}\label{eq:C10}
\]
The average is normalized by the length of the mass dyad. The indicator
denotes a fixed band covering the enhanced floor, when present. The
constants are uniform over optimizers with the stated outer bounds.
\end{lemma}
\begin{proof}
\leavevmode\par
\paragraph{Second moments and the enhanced floor.}

For the nonfloor contribution to \eqref{eq:C10}, use the whole-cap trace and bracket estimates PE (20)--(21), with cap width \(\le M_j^{-1}\), \eqref{eq:C6} giving the same \(\Delta S\) and contact hypotheses, and the admissible changed clocks giving \eqref{eq:A1}. Enlarge bands to include complete caps with comparable masses by regularity and the small radius width. This gives the first two terms. Floor mass from \(k\) uses simply \eqref{eq:C8}.

For a nonzero extra floor, use sliding radius windows on \(I\), where
\(\alpha'\ge\kappa\) and \(S\lesssim u\). Equation
\eqref{eq:A1}, integrated in radius, gives
\[
 \int\mathbb EW\,dr\lesssim x^6/C_0,
 \qquad
 \int\mathbb E\operatorname{tr}C_v^2\,dr
       \lesssim x^6/(C_0u).
\]
Here \(C_v\) denotes the spin covariance and the masses are comparable
to \(u\). Let \(\overline L^{\,+}_{r,h}\) and
\(\overline L^{\,-}_{r,h}\) be normalized averages on the following
and preceding radius windows of length \(h=t'u\). Applied in mass
coordinates, their deterministic weight is \(1/(h\alpha')\).
Let \(I_h\) omit radius strips of width \(h\) at the ends of \(I\),
so both windows remain where \(\alpha'\ge\kappa\).
Equation \eqref{eq:A1} and Fubini give
\[
 \int_{I_h}\operatorname{Var}(\overline L^{\,\pm}_{r,h})\,dr
 \lesssim\frac1{\kappa h}\int_I\mathbb EW\,dr.
\]
The single inverse power of \(h\) comes from integrating the moving
window, whose multiplicity is at most \(h\).

For two ordered cuts \(i<j\), use the identity
\[
 L_j-L_i=(y-X_j)\cdot(X_j-X_i)+\tfrac12|X_j-X_i|^2.
\]
Its noise term has second moment at most
\[
 (\mathbb E\operatorname{tr}C_j^2)^{1/2}
 \bigl(S_j-S_i+C'(D_i+D_j)\bigr).
\]
Average this identity to the right and left of the tested cut.
The radius-window kernels
\[
 K_h^\pm(i,j)=h^{-1}\mathbf1_{\{0<\pm(j-i)<h\}}
                       \mathbf1_{\{i\in I_h\}}
\]
are supported in \(I\times I\), and each marginal has density at most
one. Monotonicity and \(\operatorname{osc}_I S=O(u)\) give
\(\iint K_h^\pm(S_j-S_i)^2\,di\,dj=O(hu^2)=O(t'u^3)\).
Also \(\iint K_h^\pm(D_i+D_j)^2\le4\int_I D^2\).
Define
\[
 {\cal D}=u^{-3}\int_I D^2\,dr,\qquad
 {\cal T}=u^{-3}\int_I\mathbb E\operatorname{tr}C_{\alpha(r)}^2\,dr
       \lesssim A^3.
\]
Cauchy--Schwarz therefore bounds the integrated noise, divided by
\(u^3\), by \(C\sqrt{{\cal T}(t'+{\cal D})}\).
Combining the side-average variance bound, the raw endpoint strips,
and \(D^2\le4\operatorname{Var}L+\mathbb E\operatorname{tr}C_v^2\),
we obtain
\[
 {\cal D}\le C\left(t'+\frac{A^3}{\kappa t'}+{\cal T}
                    +\sqrt{{\cal T}(t'+{\cal D})}\right).
\]
Young's inequality absorbs the part involving \({\cal D}\):
\(C\sqrt{{\cal T}(t'+{\cal D})}
\le({\cal D}+t')/2+C'{\cal T}\).
Thus the two-sided sliding estimate is
\[
 \frac1{u^3}\int_I D^2\,dr
       \lesssim t'+A^3/(\kappa t')+A^3,\qquad 0<t'\le c.
\]
Boundary strips of width \(h\) use the raw moments. Choose \(t'\)
to balance its first two terms, or use the raw bound if that choice
exceeds the allowed interval. After multiplying by the extra density
\(\kappa\), its contribution is
\(O(A^{3/2}\sqrt\kappa)\le O(B_f^2)\). This supplies the final
term of \eqref{eq:C10}.



\end{proof}

\subsubsection{Calibration and closure under strong removal}

In addition to the outer-regularity hypotheses, suppose now that
$C(w)\gtrsim w^2$, and set
\[
 W=\max(w,\sqrt{C(w)})\asymp u\sqrt\Lambda=o(u).
\]
The endpoint estimates \eqref{eq:P9} are assumed on the regular outer
region. This stronger setting is used only in the strong-removal branch
of the profile contradiction.

\begin{lemma}[A common calibration cut]\label{bud:common-star}
On compact target bands with $cw\le v\le Cw$ and $r=O(w)$,
\[
 \int D^2\,dv=o(w^3),\qquad \int D^2\,dr=o(w^3),
 \qquad \int_0^{C'W}D_v\,dv=o(w^2)
\]
for every fixed $C'$. There is a deterministic nonfloor cut $*$ in a
regular band of scale $W$, above any prescribed fixed multiple of $w$,
for which $S_*\asymp W$ and
\[
 D_*\le x^{c_*}w^2/W,\qquad
 \|Q_*-S_*\|_p\le C_p x^{c_*}W
\]
for every fixed $p$ and some $c_*>0$ independent of $p$.
The cut is chosen from deterministic second moments and geometry before
spin sampling; it is the same cut for all these moment estimates.
\end{lemma}
\begin{proof}
\leavevmode\par
\paragraph{Inner mass and radius estimates.}

For \(v\gtrsim W\) the right side of \eqref{eq:C10} is polynomially smaller than \(w^2/W\) (\(v C(v)\gtrsim W^3\)); for the extra floor this uses
\[
 u B_f W/w_c^2\lesssim\lambda^{s_0/2+k_c/12-0.24}.
\]
On regular bands between a large fixed multiple of \(w\) and scale
\(W\), \eqref{eq:C10} saves a power relative to
\(w^2/(W^2v)^{1/3}\). The extra floor is absent there.
We also need
\[
 \int D^2=o(w^3)
\]
on bands with mass between fixed positive multiples of \(w\) and
radius \(O(w)\), separately in mass and radius. In radius, apply the
sliding estimate with \(S=O(w)\), field slope \(C_0\), and floor
lower bound \(k(O(w))\). Its errors vanish because
\(mC_0k(O(w))w^4\gg1\); narrow domains and boundary strips use
\eqref{eq:C8}.

In mass, the floor contributes \(o(w^3)\), since its mass is \(o(w)\)
and the raw scale is \(O(w)\). On caps and contact, use PE (20)--(21)
with that same raw scale. If a cap crosses the lower cutoff \(cw\)
from below \(cw/2\), use its portion above \(cw/2\). This portion
is a fixed fraction of its total mass, because outer comparability
puts its upper endpoint at \(O(w)\). The whole-cap bound on
\(\Delta S\) and the constant cap density give the trace estimate
on that portion. Thus the trace bound is always used with a positive
lower mass cutoff. The condition \(mC_0w^4\gg1\) finishes the
mass estimate. Summing these bands and the regular intermediate
estimates, and using \eqref{eq:C8} below them, gives
\(\int_0^{C'W}D\,dv=o(w^2)\) for every fixed \(C'\).


\paragraph{A good calibration cut.}

Choose a regular band of scale \(W\), above any required fixed multiple
of \(w\). The constant removal, \eqref{eq:C6}, and \eqref{eq:C10}
give a nonfloor cut \(*\) there with \(S_*\asymp W\) and
\(D_*\) polynomially smaller than \(w^2/W\). This selection uses
deterministic moments of the fixed optimizer and parameters. It is made
before sampling spins and independently of a moment order.

At this same cut, \(\|Q_*-S_*\|_p\ll W\) for every fixed \(p\),
with a saving exponent independent of \(p\). To prove this, take earlier
and later radius buffers of length \(h'=x^{c_\theta}W\), with fixed
\(c_\theta>0\) sufficiently small. Integrating \eqref{eq:C6} to or
from \(*\) shows that \(S\) varies by \(O(h')\) on fixed
enlargements: use \({\cal H}(*)\le0\), the local lower bound
\({\cal H}\ge-O(C_0/M_j^2)\), and monotonicity of \(S\).
Equations \eqref{eq:C9} and \eqref{eq:A1} make the centered averages
on these buffers polynomially smaller than \(W\). The past/future brackets and common-past comparison in
Lemma~\ref{global:buffer-brackets} then give the same improvement for
\(\||X|^2-S_*\|_p\) and for squared path increments in \(L^p\).
Finally project both leaves to a central cut a fixed fraction of
\(h'\) beyond their fork. Equations \eqref{eq:A3} and \eqref{eq:C8},
with synchronized residual layers, bound its residual by
\(O_{p,\circ}(x^3\sqrt{W/h'}/(W\sqrt{C_0}))\).
The two normalized errors just used are bounded, up to subpower losses,
by
\[
 [mC_0k(W)W^4]^{-1/2}(W/h'),\qquad
 [mC_0W^4]^{-1/2}(W/h')^{1/2},
\]
respectively. Both save a power when \(c_\theta\) is small enough,
independently of the fixed moment order. This radius-window argument
uses the floor lower bound \(mC(v)k(v)v^4\gg1\); it does not
require the upper-cap bound used in PE's original mass-window proof.



\end{proof}

Set
\[
 \widehat K=h+C(r)r+(b-C(r))S,
\]
and let \(\Gamma_{\widehat K}\) be the scalar moment curve for this
entire clock path.

\begin{proposition}[Closure for the modified constraints]\label{bud:modified-closure}
Along any fixed sequence of optimizers with the strong-removal hypotheses
above and fixed outer-comparability constants, the following estimate holds
for targets \(i\) with
\(cw\le v_i\le Cw\) and \(r_i\le Cw\),
\begin{equation}\label{eq:bud-coarse-closure}
 \frac{|S_i-\Gamma_{\widehat K}(\widehat K_i)|}{w^3}
 \le C'\{\varepsilon_m+(D_i/w)^\vartheta\},
 \qquad \varepsilon_m\longrightarrow0,\quad\vartheta>0.
\end{equation}
The constants and modulus are common to the fixed target band. Thus
this difference is \(O(w^3)\) throughout the band, and is
\(o(w^3)\) along any moving target with \(D_i=o(w)\).
Since the target-band integrals of \(D^2\) are \(o(w^3)\), the same
modulus gives, separately in the two coordinates,
\begin{equation}\label{eq:bud-coarse-closure-integrals}
 \frac1w\int
 \frac{|S_i-\Gamma_{\widehat K}(\widehat K_i)|}{w^3}\,dv_i=o(1),
 \qquad
 \frac1w\int
 \frac{|S_{\alpha(r)}-\Gamma_{\widehat K}(\widehat K_{\alpha(r)})|}
      {w^3}\,dr=o(1).
\end{equation}
Each integral is restricted to the target band just described.
\end{proposition}

\begin{proof}
The constraints \eqref{eq:C5} differ from those in the statement of
PE's closure proposition. We therefore use its proof, and verify the
replacement inputs below. A \emph{block} here consists of leaves sharing
through a fixed continuity cut of mass comparable to \(W\). A boundary
edge joins two such blocks. The tests used in that proof are the fixed finite products of centered
overlap factors and projected factors in \cite[Definition~4.2]{PE};
their natural raw size is denoted by \(p_0\).
The constants may depend on that fixed degree. All centers are the
actual moments \(S\), held fixed during removal of one site.

\begin{enumerate}[label=(\roman*)]
\item \emph{Projection estimates and row changes.}
Use the selected star \(*\) and a regular continuity cut at a large
fixed multiple of \(W\) to form the blocks. For each projection at
scale \(v\), outer comparability supplies a preceding radius interval
of length comparable to \(v\), after the low fork, with masses
comparable to \(v\). Its field length is at least
\(C(v)v/O(\log m)\). Above a sufficiently large fixed multiple of
\(W\), \(vC(v)\gtrsim W^3\). Define
\(z_v=w^3/(W^{3/2}\sqrt v)\). The postprojection bound and the
star-anchor version of \eqref{eq:A3} then read
\[
 \frac{x^3}{v\sqrt{C(v)}}
 \lesssim (x/w)^3 z_v,
 \qquad
 \frac{x^3\sqrt W}{\sqrt{v^3C(v)}}
 \lesssim (x/w)^3\frac{w^3}{Wv}.
\]
The factor \((x/w)^3\) is power small. These are the two boundary
estimates in PE (13), (16), and its Section~4, ``Analytic input and
finite-degree tests''. They also hold for conditional replacements.
The raw bounds \eqref{eq:C8}, the good-star bound, and the second-moment
estimates above supply its remaining inputs (12), (17)--(19), with
\(m,w,W\) in place of \(n,x,W\).

The clock \(\widehat K\) is monotone, bounded, and \(O(v+w)\).
Above the regularity threshold this follows from outer comparability,
\eqref{eq:C8}, and \eqref{eq:P9}; below it use monotonicity.
Consequently the scalar row sum \(\chi\) in that proof is bounded
above and below by positive constants, and its two-spin coefficient
at distinct high mass \(v\) is \(O(v)\).
The bounded-row comparison of \cite[Lemma~4.3]{PE} applies: all
clocks remain bounded and genuine, and the ordinary field increments
on the remaining sites do not change. Their raw projection norms also
follow by conditional Jensen from positive overlap tests.

\item \emph{The finite-degree expansion and calibration.}
The row interpolation replaces the cavity-spin kernel by
\[
 h+C(r)r+(b-C(r))Q^{-},
\]
where \(Q^{-}\) is the overlap with the deleted site omitted but with
the original normalization. Its difference from \(\widehat K\)
is \((b-C(r))(Q^{-}-S)\), up to the \(O(m^{-1})\) common-noise
convention. Every row interpolation preserves positive covariance.
Thus the row expansion in \cite[Lemmas~4.5, 4.7 and 4.8]{PE},
(23)--(25), has the same algebra and the same centers.
Here is how its errors are bounded after substitution.

A high error integrated with bounded split density, including repeated
vertices in one dyad, costs a power-small multiple of \(w^2/W\).
The error at \(*\) has the same bound. Ordinary low factors cost
\(O_p(W)\), or \(O_p(w)\) if their split is at the target scale.
A varying low error costs \(\int_0^{C'W}D=o(w^2)\); an error
at the old target costs \(D_i+O(m^{-1})\), with positive powers of
\(D_i/w\) after moment interpolation. Deleting the row costs
\((mw)^{-1}=o(w^3)\). The observed-block truncation step uses
variations power-smaller than \(Wz_v\) and \(p_0z_v/w\), which
were supplied in (i). These estimates justify the truncation before
any signed sums are bounded absolutely.

For completeness, suppress the common raw test factor \(p_0\).
In the terms with three low insertions, an integrated error is
\(o(w^2)\) and its other low factors cost \(O(W^2)\) or
\(O(W^3)\). If the error is at the old target and \(w\ll W\),
one additional factor has scale \(w\). When necessary, this factor
is obtained by turning off revelation below a smaller block cut of
scale \(w\), then restoring it once; the zero-field sign symmetry
cancels the term before restoration. High remainders contain either
one high error and two low factors, or two high errors and one low
factor. Their bounds are power-small multiples of \(Ww^2\) and
\(w^4/W\), respectively. Moments slightly above two suffice for
the integrated errors; their interpolation loss is paid by the strict
power saving.

If a third block is already observed, first project the fresh endpoint
in the other target block to the star and collapse its scalar row to
\(\chi\). Next replace the factor at the second leaf of the third
block by its value at the first leaf. The postprojection estimate in
(i) gives a power-saved \(w^3/(Wv)\) for this replacement. Its
scalar coefficient is \(O(v)\), and the other low factor costs
\(O(W)\), so the total is \(o(w^3)\le o(Ww^2)\). After this
replacement that leaf's
scalar row also collapses to \(\chi\). These are the two row
reductions in \cite[Lemma~4.7, observed-third-block case]{PE}.

Let \(v_{\rm blk}\asymp W\) denote the continuity cut defining the
blocks. The tests have no overlap factor internal to an old block;
its old internal relationships occur only in fixed scalar coefficients
and the prescribed topology.
Only the unweighted attachment of the first leaf in the third
block remains. At its fresh-at-cut value the summand is independent of
the deeper attachment, so the signed row-sum rule, including all old-fork atoms and coincidences,
gives the factor
\(v_{\rm blk}=O(W)\). At the star the baseline product difference
is \(O(WD_*)\); its contribution is therefore
\[
 O(v_{\rm blk}WD_*)=O(W^2D_*)=o(Ww^2).
\]
Restoring the deeper attachment uses the boundary comparison of
\cite[Lemma~4.5]{PE}; its high-layer bound is
\[
 \varepsilon_m v(Wz_v)(z_v/w)
 =\varepsilon_m w^5/W^2\le\varepsilon_m Ww^2,
 \qquad \varepsilon_m\to0,
\]
where \(\varepsilon_m\) has the fixed power saving supplied by (i). If the third block is unobserved,
its varying low error costs \(W\int_0^{O(W)}D=o(Ww^2)\);
a star attachment has the same entry factor \(O(W)\).
At the old target the two low factors have scale \(w\), with the
positive-power modulus in \(D_i/w\) described above. Thus both
observed and unobserved third blocks obey the required remainder bound.


For the leading scalar-row replacement, we spell out the test
difference in \cite[Lemma~4.8, proof of the first replacement]{PE}.
By linearity take an admissible monomial
\[
 P_a=H Y_a^k E_{oa},\qquad P_c=H Y_c^k E_{oc},\qquad
 Y_d=X_*\cdot y_d-S_*,\qquad E_{od}=Q_{od}-S_j,
\]
where \(H\) uses only exterior leaves, the incoming edge has split
\(j\in\{i,*\}\), and the edge factor is omitted if absent.
The leaf \(c\) is a fresh relative of \(a\) in its block.
Let \(\mathcal F\) be the complete target-trunk prefix through
\(*\), excluding exterior continuations and subsequent target
innovations. Define \(\Delta P_{ac}\) by averaging \(P_a-P_c\)
over the entire exterior joint continuation, conditional on
\(\mathcal F\) and the joint future of \(a,c\).
The signed exchange identity is
\[
 \sum_c\gamma_{ac}\nu[(Y_c-Y_a)P_a]
 =\frac12\sum_c\gamma_{ac}\nu[(Y_c-Y_a)\Delta P_{ac}],
\]
where \(\gamma_{ac}\) is the scalar two-spin coefficient with
revelation below the block cut moved to that cut, and the sum uses the
signed attachment rule within the block. This is the source's first
symmetrization, followed by conditional expectation.

Put \(\mathfrak s_i=w\), \(\mathfrak s_*=W\), and let \(p_H\) be the raw scale
of \(H\). Then \(p_0=p_HW^k\mathfrak s_j\) with the incoming edge
and \(p_0=p_HW^k\) without it. The exterior anchor
\(V=\mathbb E[Hy_o\mid\mathcal F]\) obeys
\[
 \|V\|_p\le\|H\|_{2p}\|Q_j\|_p^{1/2}
             \lesssim_p p_H\sqrt{\mathfrak s_j}.
\]
This is the conditional Cauchy--Schwarz estimate in that same source
proof. It uses the ordinary split-\(j\) marginal; no exterior
terminal spin is conditioned upon. The target projection powers stay
outside this averaging and outside the anchor.
For a regular innovation interval before a high split of scale \(v\),
write the unit-anchor bound as
\[
 \pi_v=O_{p,\circ}\left(\frac{x^3}{\sqrt{v^3C(v)}}\right)
 \le\varepsilon_m\frac{w^3}{W^{3/2}v},
\]
with \(\varepsilon_m\) power small. Then
\(\|Y_c-Y_a\|_p\lesssim\sqrt W\pi_v\).
Telescoping \(P_a-P_c\), a changed projection power costs
\(p_0\pi_v/\sqrt W\); the changed incoming edge, after the exterior
average, costs \(p_0\pi_v/\sqrt{\mathfrak s_j}\). Apply projection before
H\"older with the remaining powers, as in PE's test-difference
estimate. Thus
\[
 \|\Delta P_{ac}\|_p
 \lesssim\frac{\varepsilon_m p_0}{v}
       \left(\frac{w^3}{W^2}+\frac{w^{5/2}}{W^{3/2}}\right).
\]
The norm is of the exterior-averaged difference, not of the original
incoming-overlap difference. At a distinct split, the attachment
variation per dyad and \(\gamma_{ac}\) each contribute \(O(v)\).
Consequently the symmetrized error per dyad is at most
\[
 C v^2\left(\frac{\varepsilon_m w^3}{Wv}\right)
 \frac{\varepsilon_m p_0}{v}
       \left(\frac{w^3}{W^2}+\frac{w^{5/2}}{W^{3/2}}\right)
 \lesssim\varepsilon_m^2p_0Ww^2,
\]
since \(w\le W\). The saving pays the dyadic sum.
For a connected high insertion, one \(O(v)\) placement factor
suffices, because its integrated high error supplies the additional
saved \(w^2/W\); this remains true when the vertices share a dyad.

This verifies the remainder estimates in PE (24)--(25).
No variation of the optimizer occurs in this finite-tree argument.
We spell out how the finite recursion preserves their uniform target
modulus. For a block $\mathscr A$ at the chosen continuity cut
$v_{\rm blk}$, let $\nu^{\rm sc}$ be the scalar law with revelation
below $v_{\rm blk}$ moved to that cut, and
put $\gamma_{ac}=\nu^{\rm sc}[\tau_a\tau_c]$. The connected insertion
of PE (23), paired with a test $P$, is
\[
 \mathcal L(P,a)=\frac12
 \sum_{\substack{c,e,f\in\mathscr A\\ e\ne f}}
 \bigl(\nu^{\rm sc}[\tau_a\tau_c\tau_e\tau_f]
             -\gamma_{ac}\gamma_{ef}\bigr)
 (b-C)_{ef}\,\nu[P(Q_{ef}-S_{ef})].
\]
Here $S_{ef}$ is the actual moment at the split of $e,f$, and all
three added labels use the signed within-block attachment rule.
The coefficients and expectations are evaluated on the same final
topology. This formula defines the insertion through its pairing;
no positive probability interpretation of the signed sum is used.

At the star, $\widehat K_*=T_mS_*+O(W^3)$ and the changed matrix is
$T_m+O(W^2)$. Put $J_0=T_m\chi^2-1$. The recursion of
\cite[Lemma~4.9]{PE} divides by $\chi\widehat K_*\asymp W$.
The remaining pairing of the same degree at the partner root is removed
by applying the row identity to two fresh interchangeable roots absent
from $P$. Their equal leading pairings have total coefficient
$2\chi\widehat K_*$; every remaining term adds a star factor at cost
$O(W^{-1})$.
At depth $k$, accumulated coefficients are $O_k(W^{-k})$ and every
unresolved term has $k$ new centered star factors. Lemma~\ref{bud:common-star}
therefore bounds its terminal contribution by $O_k(p_0x^{kc_*})$.
Choose $k$ fixed so that $kc_*>2$, using $w\ge w_c\gg x$;
then this is $o(w^2p_0)$. Only after choosing $k$ do we fix the finite
degrees and H\"older orders. The intermediate errors give
\[
 |\mathcal L(P,a)|\le Cw^2p_0
 \left\{\varepsilon_m+(D_i/w)^\vartheta
                    +\varepsilon_m|J_0|/W^2\right\},
 \qquad \varepsilon_m\to0,
\]
with the target term absent for star-only tests. There are only finitely
many target exponents, so their positive minimum gives one $\vartheta$
for the whole fixed target band.

Return to the finite-degree row identity with $P=1$, as in
\cite[Lemma~4.10]{PE}. Its surviving high corrections are pairings
against one star factor, of raw size $p_0=W$. Apply the preceding
bound to those pairings and use
$\Gamma_{\widehat K}(\widehat K_*)=\chi^2\widehat K_*+O(W^3)$.
Indeed, the second scalar coefficient is a two-edge path through an
unobserved third block: its $O(W)$ entry factor multiplies two low
factors of size $O(W)$, and bounded finite-label variation controls
the third-order remainder.
This gives
\[
 S_*=\chi^2T_mS_*+O(W^3)
       +O\!\left(\varepsilon_m Ww^2(1+|J_0|/W^2)\right).
\]
Dividing by $S_*\asymp W$ gives
\[
 |J_0|\le CW^2+\varepsilon_m w^2
                    +\varepsilon_m|J_0|w^2/W^2.
\]
Since $w\le W$, the last term is absorbed. Hence $J_0=O(W^2)$,
and the calibrated bound is
\[
 |\mathcal L(P,a)|\le Cw^2p_0
                \{\varepsilon_m+(D_i/w)^\vartheta\}.
\]
The common base $T_m$ has remained fixed throughout the finite tree;
no rate for $T_m-T$ enters these estimates.

\item \emph{Return from the calibration scale to the target.}
Repeat the mean expansion with the block cut at a large fixed multiple
of \(w\). The two regular-dyad bounds make high errors power-smaller
than \(w\), and all low factors now have raw scale \(w\).
An old-target error contributes a positive power of \(D_i/w\),
while the remaining errors tend to zero uniformly over the target band.
The surviving mean equation at this smaller block cut is
\[
 S_i-\Gamma_{\widehat K}(\widehat K_i)
 =2(b_i-C(r_i))\chi_w\,
       \mathcal L_w(Y_a^{(i)},a)+\operatorname{Rem}_i,
 \qquad Y_a^{(i)}=X_i\cdot y_a-S_i.
\]
Here $\mathcal L_w$ is the same connected pairing with block cut of
scale $w$, and $\chi_w$ its scalar row sum. The remainder is bounded
by $Cw^3\{\varepsilon_m+(D_i/w)^\vartheta\}$, uniformly in the
fixed target band. The changed matrix coefficient is $b_i-C(r_i)$.
To use the calibrated insertion, first remove insertion edges between
this cut and the calibration block cut $v_{\rm blk}$. Their cost is
\[
 O(w)\int_0^{C'W}D=o(w^3).
\]
Next change the high scalar kernels; above a sufficiently large fixed
\(c'W\), the cost is
\[
 O(wW)\int_{c'W}^1D=o(w^3).
\]
These are precisely the two cutoff changes in PE's ``Return to the
target cut''. Allocate the error pair first relative to the sole old
root $a$ in its block, and then the summed leaf $c$; this gives the
bounded split density used in these integrals without a high-fork
atom. After the scalar-kernel change, independent block flips force
$c$ into $a$'s block, while the error endpoints share a block; if that
block is different, the connected scalar covariance vanishes.
The remaining pairing is therefore $\mathcal L(Y_a^{(i)},a)$ at the
calibration-scale block cut $v_{\rm blk}$. The test $Y_a^{(i)}$ has raw scale $p_0=w$ and is admissible by
its fresh-branch representation. Its calibrated bound is therefore
$Cw^3\{\varepsilon_m+(D_i/w)^\vartheta\}$; $b-C$ and $\chi_w$
are bounded. This proves \eqref{eq:bud-coarse-closure}. The already proved target-band
\(D^2\) bounds then give both integrals in
\eqref{eq:bud-coarse-closure-integrals}.
\end{enumerate}

For example, when \(0<\vartheta\le2\), H\"older bounds its discrepancy
term by a constant times
\( (w^{-3}\int D^2)^{\vartheta/2}\); a larger exponent may first
be decreased using the raw bound \(D=O(w)\).

\end{proof}

\begin{proof}[Proof of Proposition~\ref{bud:coarse-profile}]
We first exclude a failure of comparability, using a PE child for weak
removal and Proposition~\ref{bud:modified-closure} for strong removal.
The final step identifies the initial slope by a macroscopic transfer.

Profile arguments allow subselection. If loose fixed comparability
fails, choose an offending mass $w\ge w_c$ comparable to the supremum
of offending masses. This supplies the outer regularity above a fixed
multiple of $w$. Choose the loose constants strictly outside the bounds
of the PE direct profile limits. The scale estimates displayed before
Lemma~\ref{bud:outer-moments} follow from $w\ge w_c$.

\paragraph{Step 1: weak removal and feasible barriers.}

First suppose \(C(w)=o(w^2)\). Apply PE finite propagation to a child acting from the changed path of this optimizer, with the data \(d',E'\) of \cite[Lemma~6.2]{PE} at \(w\), and ordinary PE constraints. The data verification in that lemma uses (3) and the parent strength, cutoff and outer bound (not the parent's optimization constraints); hence (3)--(4) hold in the child. It has nonnegative monotone bare paths with bounded clocks and a fixed positive all-mass matrix margin, choosing its fixed \(d'/w\) small. Its final error is at most \(m^{20\delta}E'\); write \(z\) for a final minimizing quantile and \(\mathcal D(z)=d'^2 f(z/d')\). The movement inequality of \cite[Lemma~6.3]{PE} applies with parent-removal \(C\) to any trial satisfying \eqref{eq:C5}: its proof only uses feasibility, optimality, the slope bound on \(C\) and the square interpolation. We explain the corresponding barriers to handle the modified floor.
Here \(d'=c_0 w,\ E'=\min(m^{-1/2}d'^2,m^{-100\delta}C(w)w^3)\), \(c_0>0\) sufficiently small fixed. In particular the parent changed endpoint still tests against the common \(T_m\) by the estimates of \cite[Lemma~6.2]{PE}, since \(r=O(y)\) on comparable mass bands of scale \(y\gtrsim w\), \(r=O(w)\) below that.

Suppose that \(r\) and \(z\) are separated by fixed thresholds whose
gap is comparable to \(y\), on a mass interval of length comparable
to \(y\), where \(y\gtrsim w\). Macro intervals lie strictly inside
\((0,1)\). Failure of agreement in normalized measure produces such
an interval by monotonicity, and outer regularity gives \(r\le C'y\)
there. Here agreement means that, on each fixed comparable mass band,
the mass of \(\{|r-z|>\epsilon y\}\), divided by \(y\), tends to
zero for every fixed \(\epsilon>0\).

The forward constraints are
\(l\le r'\le1/\bar k(r)\), with \(l=M_j^{-1}\). We have
\(\sup\bar k=o(1)\) and \(\int_0^1\bar k=o(w)\), the latter
using \(w\gtrsim u\sqrt\Lambda\). Reserve fixed short portions
at the two ends of the separation interval for the following barriers.
If \(z\) is above \(r\), choose an intermediate threshold \(H\)
and an early subinterval \([a,b]\), with \(b-a\asymp y\), and set
\[
 b_\uparrow(v)=lv+
 \begin{cases}
 0,&v\le a,\\
 H(v-a)/(b-a),&a<v<b,\\
 H,&v\ge b.
 \end{cases}
\]
Choose \(H\asymp y\) in the threshold gap, leaving fixed distance
from its ends. The rise has fixed slope
\(s=l+H/(b-a)\), and eventually \(s\sup\bar k\le1\).
Also \(l+H<1\), so \(\max(r,b_\uparrow)\) preserves both
endpoints and both slope bounds.

If \(z\) is below \(r\), use \(b_\downarrow(v)=H+lv\) until
a late point \(a\). Thereafter solve
\[
 \int_{H+la}^{b_\downarrow(v)}\bar k(t)\,dt=v-a
\]
until this curve meets \(1-(1-v)l\), and follow that line after the
meeting. The ramp has derivative \(1/\bar k\). Since
\(l\sup\bar k<1/2\), its mass length is at most
\(2\int_0^1\bar k=o(w)\), so it fits in the reserved part of the
interval. The trial \(\min(r,b_\downarrow)\) again preserves both
endpoints and both bounds. In both constructions choose the threshold
and reserved portions as in \cite[Lemma~6.4]{PE}, so every moved point
stays distance \(\gtrsim y\) from \(z\). Upward motion has radii
\(O(y)\), with \(r\asymp y\) if \(y/w\to\infty\); downward
motion has both parent radii comparable to \(y\).

If \(w\ge\lambda^{k_c/4}u\), full clipping already satisfies the strict movement fraction (displaced distance less than \(|z-r|\mathcal D(z)/(\mathcal D(z)+C(r))\)). For \(y/w\to\infty\), \(y\mathcal D(y)\gtrsim w^3\gg y C(y)\) since \(y C(y)\lesssim_\circ u^3\Lambda\); at \(y\asymp w\), \(C(0)\ll w^2\). Thus for \(z\le O(y)\) the ratio of strengths is sufficient with our placement margin, and for \(z/y\to\infty\) use monotonicity of \(z\mathcal D(z)\) to permit any fixed \(O(y)\) displacement.

In the remaining case \(w<\lambda^{k_c/4}u\) we need separation tests only up to \(y\le m^{2\delta}w\ll u\), since the child is already comparable above \(m^\delta d'\) by finite propagation. All active radii are now below the enhanced patch. Use the partial interpolation of \cite[Lemma~6.4]{PE} with fraction a small constant times \(\min(1,\mathcal D(y)/C(y))\), forward for upwards, affine in \(\int_0^r k\) for downwards. On the active portions \(\bar k=k\); feasibility and the strict fraction proof in that lemma use just \(l\) as lower slope, convexity of the PE \(k\) in its integrated coordinate, \(C(r)\asymp C(y)\) for \(r\asymp y\), \(C(0)=o(w^2)\) at \(y\asymp w\), and monotonicity of \(z\mathcal D(z)\). Hence they are unchanged.

For either trial write \(\widehat r=r+\omega(z-r)\). The movement
inequality of \cite[Lemma~6.3]{PE} states
\[
 \int C(r)\omega(z-r)^2\,dv\le4m^{20\delta}E',
 \qquad 0\le\omega<\frac{\mathcal D(z)}{C(r)+\mathcal D(z)}.
\]
The constructions just verified make its left side at least
\(c y^3\min(C(y),\mathcal D(y))\gg m^{20\delta}E'\).
Thus no fixed threshold separation can occur. The resulting agreement
in normalized measure gives the child outer regularity above \(w\)
when \(w=o(1)\), and agreement on compact positive bands of scale
\(w\). Its direct-test hypotheses hold there with
\(\mathcal D(w)\asymp w^2\) and \(mE'^2\lesssim w^4\).
Apply \cite[Theorem~5.1 and Proposition~5.9]{PE} to the child.
Its continuous limiting profile transfers to the parent pointwise:
bracket any tested mass between two nearby continuity masses, use
agreement in measure to select comparison points on either side,
and then use monotonicity of \(r,z\). Shrinking the bracket gives
the child's loose bounds also at the offending point. Near mass one
use \(r\le1\) as well. This excludes a weak-removal violation.




\paragraph{Step 2: strong removal and the profile contradiction.}

Suppose $C(w)\gtrsim w^2$. Then
$w\le W=\max(w,\sqrt{C(w)})\asymp u\sqrt\Lambda=o(u)$,
so the conditional analytic hypotheses hold. The bare endpoints satisfy
\eqref{eq:P9} on the regular exterior. Lemma~\ref{bud:common-star} and
Proposition~\ref{bud:modified-closure} therefore provide the common
calibration cut and the closure in both coordinates required next.

We use this closure to identify the scaled optimizer profile,
following PE (28)--(29). On each compact band in the scaled radius,
\(p=0\) and \(C=C_0\). Cap widths remain negligible after
multiplication by \(C_0/w^3\), and the floor mass on each such band
is \(o(w)\). Equation \eqref{eq:C6} has no additional
state-dependent factor there.

Subselect the distribution of \(r/w\) under mass divided by \(w\).
It converges locally to a Radon measure on \([0,\infty)\), whose
cumulative function we denote by \(U_1\). Outer regularity gives
\(U_1(t)\asymp t\) for large \(t\). The scalar clock
\(\widehat K/(T_mw)\) has the same limiting distribution.
Indeed on compact positive bands of scaled mass, \(r,S=O(w)\),
\(|S-r|\le M_j^{-1}\) almost everywhere off the negligible floor,
and \(((b-T_m)S+h)/w^3=o(1)\). Very small scaled masses contribute
at most their length. Monotonicity and the comparison on large fixed
bands exclude far scaled masses from any fixed clock band.

\paragraph{The profile contradiction.}

Set
\[
 g_m(t)=\frac{\Gamma_{\widehat K}(T_mwt)-wt}{w^3},
\]
extending the scalar clock if necessary. Let \(a_2,a_4\) be
subsequential limits of its second and fourth spatial continuation
derivatives at zero time and field. The scalar derivative formula and
\cite[Lemma~5.3]{PE}, applied to the bounded normalized clocks, give
\[
 g_m(0)=0,\qquad g_m''\longrightarrow At-BU_1(t)
       \quad\text{in }L^1_{\rm loc},\qquad
 A=T^3a_4^2,\quad B=2T^2a_2^3.
\]
Both constants have fixed positive lower and upper bounds, and the
second derivatives are locally bounded. To see the limit, oddness and
bounded time and space derivatives give
\(V_{zzz}(T_mwt,Z)=a_{4,m}Z+O_p(w^{3/2})\), while
\(\mathbb EZ^2/w\to Tt\). The rescaled clock mass converges at
continuity points of \(U_1\).

Choose a nonfloor point of regular scale comparable to \(w\).
Such points exist because the local floor mass is small. There
\(S\asymp w\) and \(\widehat K-T_mS=o(w^3)\), since
\(C_0|S-r|/w^3=o(1)\). Equation
\eqref{eq:bud-coarse-closure} bounds \(g_m\) at an argument bounded
above and below by positive constants. Together with \(g_m(0)=0\)
and the curvature bound, this bounds its initial slope. Pass to a
\(C^1_{\rm loc}\) limit \(g\). The mass estimate in
\eqref{eq:bud-coarse-closure-integrals} allows nonfloor approximants
to every support point, giving \(g=0\) on the support; at zero use
\(g(0)=0\).

Above the support infimum, subselect further so that
\(S_{\alpha(wt)}/w\to s(t)\) almost everywhere. On compact
intervals there, \(\alpha(wt)\asymp w\). The closure identity is
\[
 \frac{C_0(S_{\alpha(wt)}-wt)}{T_mw^3}
 =g_m\bigl(\widehat K(\alpha(wt))/(T_mw)\bigr)+\mathrm{err}.
\]
The error is bounded and tends to zero in local \(L^1(dt)\), by the
radius estimate in \eqref{eq:bud-coarse-closure-integrals} and the
bare-clock bound. If \(C_0/w^2\to h_1\in(0,\infty)\), it gives
\[
 g(s(t))=h_1(s(t)-t)/T,\qquad
 \mathcal H'(wt)/w^3\longrightarrow2h_1(t-s(t)).
\]
If \(C_0/w^2\to\infty\), it instead gives
\(S_{\alpha(wt)}=wt+o(w)\) and
\(\widehat K/(T_mw)=t+o(1)\); the derivative limit is then
\(-2Tg(t)\). These derivative sequences are locally bounded and
converge in local \(L^1\).

On an interior support gap, the potential \(\mathcal H(wt)/w^4\)
therefore has a nonnegative limiting primitive that vanishes at both
ends. Nonfloor approximants pin the endpoint values to zero; for a
negative cap, use its nearest right zero within \(M_j^{-1}\).
The same cap-width bound excludes a negative excursion inside the gap.
If the left endpoint is a first atom, choose approximants after a fixed
positive amount of scaled mass has accumulated. The derivative bound
then holds between them and the gap, even when the atom is at zero.
All these estimates are used only where scaled radii are bounded and
scaled masses are bounded above and below.

Monotonicity gives \(s(t)\) in the closed gap interval almost
everywhere. A positive derivative of \(g\) at the left endpoint
would force the limiting primitive to decrease immediately to its
right. For finite \(h_1\), the identity
\(g(s(t))=h_1(s(t)-t)/T\) forces \(s(t)>t\), whether \(s(t)\)
is near that endpoint or farther into the gap. For infinite strength,
the conclusion follows from the derivative \(-2Tg(t)\).
Both contradict nonnegativity and the zero endpoint value. Reflection
excludes a positive derivative at the right endpoint as well.
But on a gap of length \(L\), the curvature is
\(g''(t)=At-\mathrm{const}\) and the endpoint values are zero,
so the sum of endpoint derivatives is \(AL^2/6>0\).
No gap is possible.

Above the support infimum \(l\), we thus have \(g=0\) and
\(U_1(t)=At/B\). If \(l>0\), then \(g''=At\) below \(l\)
and \(g(l)=g'(l)=0\), contradicting \(g(0)=0\).
Consequently \(U_1(t)=At/B\) throughout. This is the linear quantile
limit PE (28), with the same fixed slope bounds. It contradicts our
choice of loose comparability constants and proves the first part of
\eqref{eq:C7}.

The first part of \eqref{eq:C7} has now been proved for every
optimizer. In particular Lemmas~\ref{bud:outer-moments}
and~\ref{bud:outer-second} can henceforth be used with $w=w_c$ and fixed uniform comparability constants. This is the
version of \eqref{eq:C8}--\eqref{eq:C10} used in the strengthening
argument below; it does not invoke the strong-removal closure again.

\paragraph{Step 3: identification of the initial slope.}

Repeat the weak-removal transfer at \(w=u\), now without assuming
a violation. We identify the child's linear slope by its macroscopic
clock limit. First,
\(S_{\rm child}-z\to0\) in mass, by PE (18), geometry, and
regularity on each fixed interior band. The required error estimates
are
\[
 m\mathcal D(v)\gtrsim mu^3,\qquad
 E'/d'^3\lesssim m^{-1/2}/u,\qquad
 m^{\delta/2}/(mE')\longrightarrow0.
\]
Compare \(z\) to a further ordinary PE minimizer at final strength,
with fixed small removal parameter \(d''>0\) and
\(E''=m^{-1/2}d''^2\). Its macroscopic endpoint hypothesis PE (4)
follows from \eqref{eq:P9} and the vanishing preceding removals;
all-mass admissibility and a positive covariance margin remain.
Finite propagation gives error \(m^{20\delta}E''\ll u^3\), and
\cite[Proposition~5.9]{PE} gives the further minimizer's limiting
quantile \(q_T\).

The ordinary PE barriers transfer this limit to \(z\). At a fixed
interior mass scale, full clipping is feasible with the child's minimum
slope \(1/M(v)\); both forced lengths are \(o(1)\).
The movement fraction holds because \(\sup\mathcal D=o(1)\),
whereas the newest removal is bounded below. A fixed threshold
separation would cost \(\gtrsim\inf\mathcal D\gtrsim u^3\),
contradicting the propagation error. Thus
\(h_{\rm child}+b_{\rm child}S_{\rm child}\to Tq_T\) in
macroscopic \(L^1\).

Bounded-clock PDE stability, as at \eqref{eq:G3}, now gives convergence
of the scalar spatial derivatives at zero. Equivalently one can use
drift comparison after \(L^1\) convergence of the inverse-clock
distributions. The slope formula \(A/B\) in PE (29), together with
\cite[Lemma~5.8]{PE}, therefore identifies the child's inverse-profile slope as
\(\kappa_T\). Agreement transfers its linear quantile limit to
\(r\); meanwhile \(q(v)/v\to1/\kappa_T\) by \eqref{eq:P2}
and the normalized inverse construction. This proves the second part
of \eqref{eq:C7}. On compact positive radius bands of scale \(u\)
inside the constant-removal region, it also gives \(S-r=o(u)\)
uniformly, by \eqref{eq:C6}, monotonicity, and the local \(o(u)\)
floor mass.
\end{proof}


\subsection{Finite strengthening and common parameter envelopes}

Use the notation
\[
 K_*=h_*+bq,\quad d=h-h_*,\quad K_s=h+C(r)r+(b-C(r))S,\quad
 \Gamma_*=\Gamma_{K_*},\quad \Gamma_s=\Gamma_{K_s}
\]
in the normalized local argument (\(s\) on \(K_s,\Gamma_s\) here denotes the sampled path, not a tilt). Let \(\beta_0(r)=q^{-1}(r)\) locally near the patch; it and \(q\) have slopes between fixed positive constants there. Choose finitely many nested compact intervals scaled by \(u\), strictly interior to the constant profile and at positive radii of order \(u\), each containing with room the required image of the mass band by \eqref{eq:C7}. Their fixed spacings can be as small as needed for the number of stages.

Take
\[
 a_{-1}=1,\quad
 a_j=\max(a_{\rm goal},\lambda^{-0.12} A^{1/3} a_{j-1}^{\,2/3}),\quad
 \kappa_j=(A/a_{j-1})^2\quad (j\ge0).                                   \tag{C11}\label{eq:C11}
\]
These scales are nonincreasing, \(a_0=o(1)\), and reach \(a_{\rm goal}\) in a bounded fixed number of steps: the iterated term has fixed point \(A\lambda^{-0.36}\ll a_{\rm goal}\), while \(A\gtrsim\lambda^{2/p_*}\). Use \eqref{eq:C5} with the indicated floor on the next nested patch \(I_j\), dropping the previous extra floor.

\begin{lemma}[Finite strengthening]\label{bud:strengthening}
On a slightly smaller interval interior to \(I_j\), at each stage,
\[
 |\alpha-\beta_0|\lesssim u a_j Z                                                   \tag{C12}\label{eq:C12}
\]
for almost every global parameter, simultaneously for all its optimizers, where \(Z\ge1\) has bounded averaged square in that parameter. For each fixed choice of the other path data, the exceptional set is independent of the optimizer and may be taken common to the finitely many stages. At successive stages that parameter may be restricted to slightly smaller fixed interior intervals, always leaving room beyond \([1/3,2/3]\); parameters not explicitly varied need no restriction beyond those of \eqref{eq:C3}. At stage \(-1\) the analogous assertion holds by \eqref{eq:C7} with \(Z=1\).

\end{lemma}

We prove the lemma by induction.  First we show that the preceding-stage
bound pays the new optimization and supplies a common global norm bound.
We then rule out a violation at the new stage by a buffered contradiction.
The comparison cost therefore depends only on the already proved accuracy.



\begin{lemma}[Stage cost and an envelope common to all minimizers]\label{bud:envelope}
Suppose the preceding stage satisfies \eqref{eq:C12}.  At the new stage
the average optimized excess is at most
\(O(x^{4-2\zeta}u\Lambda^{1/3})\). For almost every global parameter,
every minimizer simultaneously satisfies \eqref{eq:C13} below,
with a common measurable majorant \(Z\) whose
averaged square is bounded.
\end{lemma}
\begin{proof}

Let the old optimizer have radius law \(k\,dr+\nu\). On a fixed
enlargement of \(I_j\), convolve the fraction \(c'\kappa_j\) of
\(\nu\) with the symmetric uniform kernel of radius
\(\sigma=\min(c_{\rm small}u,C_{\rm big}u a_{j-1}Z)\).
Choose the fixed constants, depending on the stage, so that all motion
stays inside the preceding controlled band and the constant part of
\(C\). Every radius interval of length \(2\sigma\) used here carries
mass \(\gtrsim\sigma\). This follows from \eqref{eq:C12} and the
positive lower slope of \(\beta_0\); if the fixed cutoff determines
\(\sigma\), use \eqref{eq:C7} instead. Subtracting the density \(k\)
does not affect this lower bound. Taking \(c'\) sufficiently large
therefore supplies density at least \(\kappa_j\) throughout \(I_j\).
The convolved fraction is \(o(1)\), so the new upper density bound is
at most twice the old one.

Write \(r_{\rm old},r_{\rm new}\) for the two quantiles. Conditional
Jensen shows that the new radius law dominates the old in convex order
and has the same mean. For every \(a'\in[0,1]\), the upper-tail
quantile formula
\[
 \int_{a'}^1r(v)\,dv
 =\inf_{s'}\{(1-a')s'+\mathbb E(r-s')_+\}
\]
therefore gives a nonnegative upper-tail integral of
\(r_{\rm new}-r_{\rm old}\). Integration by parts against the
nondecreasing moment \(S_t\) at each forward interpolation time yields
\[
 \int_0^1 S_t(v)(r_{\rm new}(v)-r_{\rm old}(v))\,dv\ge0.
\]
Since all motion occurs where \(C=C_0\), \eqref{eq:A1} makes the
pressure change nonpositive. Only the penalty must be paid; it is
\(O(C_0\kappa_ju\sigma^2)\lesssim x^4u\Lambda^{1/3}Z^2\).
Averaging and starting from the bound before \eqref{eq:C5} gives the
asserted cost at every stage.

We next obtain a norm bound common to every new minimizer. Write
\(g=\gamma_0\), let \(F_j(g,\alpha)\) be the objective with the
inverse profile fixed, and put \(V_j(g)=\min_\alpha F_j(g,\alpha)\).
The admissible density box is independent of \(g\).
Equation~\eqref{eq:A1}, followed by its fixed-size path limit, makes
each frozen objective differentiable and bounds its derivative uniformly.
The minimum \(V_j\) is therefore Lipschitz.

At a differentiability point of \(V_j\), fix any minimizing
\(\alpha\). The function
\(g'\mapsto F_j(g',\alpha)-V_j(g')\) is nonnegative and vanishes
at \(g\). Its two one-sided derivatives give
\[
 \partial_gF_j(g,\alpha)=V_j'(g).
\]
The exceptional set is the nondifferentiability set of the single
function \(V_j\); it is independent of the minimizer.

Restore the scalar terms by setting
\[
 \mathcal R_j(g)=V_j(g)-\frac14\int bq^2\,dv
                      +\frac12\int_{v\ge x}d(v)q(v)\,dv .
\]
For each minimizer define
\[
 \mathcal N^2=\int(v^2+U^2)
                    [D_v^2+(S_v-q)^2]\,dv.
\]
The aligned matrix and field variation in \eqref{eq:P10} gives
\(\mathcal R_j'(g)\ge c\lambda\mathcal N^2\).
Above the mask this is the completed square
\(c_g[D^2+(S-q)^2]/4\); below it the field stays fixed and the
extra term is nonnegative. Thus the same derivative bounds
\(\mathcal N\) for every minimizer at \(g\).

Choose an upper parameter endpoint outside the common null set in a fixed positive-length margin
beyond the next averaging interval, where the optimized excess is at
most a fixed multiple of its average on that margin. At the lower
endpoint use the nonnegativity of the excess. Integrating
\(\mathcal R_j'\) and using the direct pressure comparison gives
\[
 c\lambda\,\mathbf E_{\gamma_0}\mathcal N^2
 \le\mathbf E_{\gamma_0}\mathcal R_j'
 \lesssim\mathcal B+x^{4-2\zeta}u\Lambda^{1/3}.
\]
This notation denotes a bound by the common measurable derivative,
not integration of a chosen optimizer. The endpoint depends only on
the stage and the fixed path data.

Choose \(0<h_e\ll\min(1,c_G)\), with \(p_*\ll h_e\).
The bound on \(\mathcal B\), with \(\zeta\) sufficiently small,
now permits the following explicit common majorant. Off the common null
set take
\[
 Z(g)^2=1+\frac{\mathcal R_j'(g)}
                    {c\lambda x^4\lambda^{-2s_0}},
\]
using the lower-bound constant \(c\) above, and set \(Z=1\) on
that null set. If \(\lambda=x^{p_*}\), \(u\le x^{h_e}\), and
\(\chi_e\le x^{h_e}\), use the smaller denominator
\(c\lambda x^{4+h_e/2}\) instead. The integrated derivative bound
gives \(\mathbf E Z^2\lesssim1\) in both cases. Thus \(Z\ge1\),
and
simultaneously for every minimizer outside that set,
\[
 \mathcal N\le x^2\lambda^{-s_0}Z,\qquad
 \mathcal N\le x^{2+h_e/4}Z
 \quad\text{if }\lambda=x^{p_*},\ u\le x^{h_e},\ \chi_e\le x^{h_e}.
 \tag{C13}\label{eq:C13}
\]
Only direct paths occur here. The majorant may change at the next
stage. Restricting to slightly smaller deterministic parameter intervals
leaves room beyond \([1/3,2/3]\), since there are only finitely many
stages. The only use of the preceding-stage accuracy was to pay for
the new floor; no estimate at the new stage was assumed.

\end{proof}



\subsubsection{Buffered geometry for a possible violation}

Fix the new stage and a sequence of nonexceptional global parameter
values contradicting \eqref{eq:C12}.
Choose a radius \(r_0\) and discrepancy
\(d_0=uL=|\alpha(r_0)-\beta_0(r_0)|\), with
\(L\ge a_jZ\) and \(L\to0\) by \eqref{eq:C7}.
We may arrange
\[
 |\alpha-\beta_0|=O(d_0)
 \quad\text{within radius }b_{\rm buf}
      =u(\sqrt L+\lambda^\rho)\text{ on either side of }r_0.
\]
Here \(\rho>0\) is fixed and sufficiently small relative to the
power gaps below. To obtain this buffer, whenever its discrepancy bound
fails, move to a point where the discrepancy has at least doubled and
repeat. There are only logarithmically many moves. The sum of the
\(u\sqrt L\) travel bounds is dominated by the last one, which is
\(o(u)\) by \eqref{eq:C7}; the sum of the
\(u\lambda^\rho\) bounds is also \(o(u)\).
Thus all selected points stay in the interior of the patch.

Put
\[
 h_1=u\max(L,\sqrt\lambda),\qquad
 G_1=\min(\max(L^2,\lambda),\Lambda),\qquad P_1=LG_1.
\]
The buffer contains unboundedly many lengths \(h_1\). On each smaller
fixed fraction, every cut has nonfloor points within \(O(d_0)\) on
both sides: otherwise the small floor density contradicts the positive
lower slope of \(\beta_0\) and the discrepancy bound.
Equation \eqref{eq:C6} and monotonicity then give
\(|S-r|=O(d_0)\), while the oscillations of \(S\) and \(\alpha\)
on a radius interval of length \(l\) are \(O(d_0+l)\).
At nonfloor points, \(|S-r|\le M_j^{-1}\).
Every later buffer restriction is by a fixed proportion; the finite
number needed in the moment iteration is specified there.

To compare scalar curvatures, define in this buffer
\[
 k_*(r)=\Gamma_*^{-1}(r),\qquad
 \widetilde r(v)=\Gamma_*(K_s(v)),\qquad
 \widetilde\alpha(r)=|\{v:\widetilde r(v)\le r\}|.
\]
Near the tested band the defect \(d\) consists only of the logarithmic
seed, so
\[
 |d|\lesssim_\circ u^3 X^3\lambda^{1+\eta-3s_*/2}=o(u^3P_1).
\]
The comparison uses \(G_1\ge\sqrt\lambda L\) and
\(L\ge X^2\lambda^{-s_0}\). Split the latter lower bound at
\(\sqrt\lambda\) and \(\sqrt\Lambda\); in each range the ratio
has a power saving. At the \(\sqrt\lambda\) breakpoint the saving
is \(1/4+O(\tau+\eta)\) in \(\lambda\).
Equation \eqref{eq:P2} and monotone bracketing now give
\(\widetilde r(\alpha(r))=r+O(d_0)\) and
\(\widetilde\alpha-\beta_0=O(d_0)\).

Extend the scalar clocks annealed to a common bounded time interval,
and write their distribution difference in mass length as
\(\delta\alpha_t=|\{K_s\le t\}|-|\{K_*\le t\}|\). We claim
\[
 \int \min(u,t)|\delta\alpha_t|dt=o(d_0).
\]
The one-dimensional inverse-transport bound, together with
\eqref{eq:C7}, \eqref{eq:C8}, and \eqref{eq:P4}--\eqref{eq:P7},
bounds the left side by
\(\lesssim\int\min(u,v+w_c)|K_s-K_*|\,dv\).
Here the low clock is controlled because
\(\lambda_1U^3\log(1/x)\lesssim w_c\).
Masses \(v\le O(w_c)\) contribute \(O(w_c^3)\ll d_0\).
Above them, the terms \(b(S-q)\) and \(d\) contribute
\(O(\sqrt u\,x^2\lambda^{-s_0}Z)\), by \eqref{eq:C13} and
\eqref{eq:P8}.

It remains to bound \(C(r)(r-S)\). Off the floor, use
\eqref{eq:C6} and the same norm. The original floor contributes
\(O(uC_0\int k)\) by boundedness. On the extra floor, at
\(r\asymp u\), the \(S-q\) part is again controlled by the norm.
Fubini and \eqref{eq:C7} bound the remaining \(|r-q|\) integral by
\[
 O(\kappa_j)\int_{I_j}|\alpha(r)-\beta_0(r)|dr
 \lesssim \kappa_j\int_{v\asymp u}|r(v)-q(v)|\,dv ,
\]
on a comparable enlargement. Its extra-floor portion can be absorbed,
since \(\kappa_j=o(1)\); its other portions have just been bounded.
Multiplication by \(uC_0\) proves the claimed transport estimate,
with a power margin for subpower losses.


\Needspace{4\baselineskip}
\begin{lemma}[Scalar curvature in a buffer]\label{bud:curvature}

By the scalar second derivative identity we now have uniformly a.e. in the inner buffer
\[
 \frac{d^2}{dr^2}\big[\Gamma_s(k_*(r))-r\big]
       =-A_* (\widetilde\alpha(r)-\beta_0(r))+o(d_0),                     \tag{C14}\label{eq:C14}
\]
where \(A_*\) converges to a positive constant (\(2(V_{*,zz}(0,0))^3 k_*'(0)^2\)). \end{lemma}
\begin{proof}
Let \(Q_0\) be a fixed bound for the extended clocks. At time
\(t\asymp u\), PDE drift comparison bounds the difference of
continuation derivatives in sup norm by
\(O(\int_t^{Q_0}|\delta\alpha_v|\,dv)\); odd derivatives gain
an additional factor \(|z|\). For a bounded smooth even observable,
changing the forward diffusion costs at most
\[
 C\int_0^t v\left(|\delta\alpha_v|
          +\int_v^{Q_0}|\delta\alpha_{v'}|\,dv'\right)\,dv.
\]
Indeed the drift difference is odd and bounded by the indicated
coefficient times \(|Z_v|\). The propagated spatial gradient is also
odd with bounded derivative, and \(\mathbb EZ_v^2=O(v)\).
These comparisons require only bounded spatial drift derivatives on
the bounded time interval.

Now use
\(\Gamma''=\mathbb EV_{zzz}^2-2\alpha_t\mathbb EV_{zz}^3\).
The two masses are \(O(u)\), and their difference is \(O(d_0)\).
Oddness in the first term and the preceding transport estimates leave,
to order \(d_0\), only the mass-distribution difference in the second
term. Multiplying by \(k_*'^2\) gives the displayed coefficient
\(A_*\). The chain-rule term involving \(k_*''=O(\lambda u)\)
uses the same comparison for \(\Gamma_s'-\Gamma_*'\) and is
\(o(d_0)\). This proves \eqref{eq:C14}.

\end{proof}




\subsection{Buffered moments, free forks, and power margins}

Write \(\delta_i=Q_i-S_i\) at a fork. For a radius interval
\(J=[a,b]\) in the inner buffer with \(|J|\asymp d_0\), define
the two box integrals by
\[
 \operatorname{avg}_{J,r}F
   =\frac1{d_0}\int_a^b F_{\alpha(r)}\,dr,
 \qquad
 \operatorname{avg}_{J,v}F
   =\frac1{d_0}\int_{\alpha(a)}^{\alpha(b)}F_v\,dv.
\]
Endpoints have zero measure in these integrals. Both use the scale
\(d_0\) as denominator. In particular the second is not normalized by
the mass of the interval: that mass is \(O(d_0)\), but can be
arbitrarily small. In box assertions, \(\operatorname{avg}\) denotes
either operation, or their sum, uniformly over moving intervals of this
kind. An explicitly subscripted dyadic average such as
\(\operatorname{avg}_{v\sim w}\) instead divides the mass integral
by the length of that dyadic band.

\begin{lemma}[Moments on moving boxes]\label{bud:box-moments}
Uniformly inside the buffer and for these averages, for every fixed
moment order, we have
\[
 \|\delta_i\|_p\lesssim_p d_0,\qquad
 \operatorname{avg}\|\delta_i\|_2^2\lesssim d_0^2\lambda^{0.14},\qquad
 \operatorname{avg}\|Q_i-|X_i|^2\|_2^2
                         \lesssim_\circ \lambda^{-\rho}u^2 A^3.          \tag{C15}\label{eq:C15}
\]

\end{lemma}\begin{proof}


Fix the optimizer and the moment order \(p\). Choose a finite sequence
of nested buffers \(B_0\supset B_1\supset\cdots\supset B_{n_{\rm buf}}\),
with successive boundaries separated by fixed fractions of
\(b_{\rm buf}\). Their number will be fixed below, before taking the
size limit. Define
\[
 {\cal Q}_k=\max\left\{1,\sup_{t,t'\,\text{in }B_k}
  \frac{\||X_t-X_{t'}|^2\|_p}
       {d_0+|r(t)-r(t')|}\right\}.
\]
The cuts in this supremum are deterministic and lie on one path.
Boundedness gives finiteness at fixed size, and \eqref{eq:C8} gives
\({\cal Q}_0=O_p(L^{-1})\). To estimate cuts in \(B_{k+1}\),
use only windows lying in \(B_k\). If such a window \(J\) has
radius length \(\Delta\gtrsim d_0\), the weighted-layer calculation
for \eqref{eq:C9} gives
\[
 \left\|\int_J W\,dr\right\|_{p/2}
 \lesssim_{p,\circ} u^3 A^3\Delta({\cal Q}_k+\lambda^{-\rho}).
\]
For layers ending within radius distance \(O(u)\) and still at
comparable scales, projection after distance \(t'\) costs
\(O_\circ(x^6/(C_0ut'))\) times the earlier squared increment,
including the Hessian weight. The normalized-anchor truncation in
Paragraph~\ref{par:global-block-estimates} retains that increment;
it requires no higher moment of \({\cal Q}_k\). When
\(t'\lesssim\Delta\), \eqref{eq:global-grid} bounds its integral by
\(O_p(t'{\cal Q}_k\Delta)\), using terminal increments on an
interval of length \(O(\Delta)\) inside \(B_k\). When
\(t'\gtrsim\Delta\), its pointwise norm is
\(O_p(t'{\cal Q}_k)\), giving the same integrated bound.
Layers leaving \(B_k\) start at distance
\(t'\gtrsim u\lambda^\rho\) and use \eqref{eq:C8}; this accounts
for \(\lambda^{-\rho}\). The first inverse-power increments are
negligible by the square-sum estimate. Outer layers at scale
\(v\gtrsim u\) cost \(O_{p,\circ}(x^6/(vC(v)))\) per unit
radius length, which also fits the displayed bound.

Using \(\alpha'\ge\kappa_j\), normalized \(L\) averages on such intervals have centered norm relative to \(\Delta\) at most
\(O_{p,\circ}((A/L)^{3/2}\kappa_j^{-1/2})({\cal Q}_k+\lambda^{-\rho})^{1/2}\).
The prefactor, before subpower losses, is at most
\(\lambda^{0.18}\) by \eqref{eq:C11}. For two cuts in
\(B_{k+1}\), take earlier and later windows of length comparable to
\(\Delta=\max(d_0,\text{their radius separation})\).
The brackets in Lemma~\ref{global:buffer-brackets} then control
both centered square norms and their squared increment. The
deterministic oscillation of \(S\) costs \(O(\Delta)\). After
absorbing subpower losses and the \(\lambda^{-\rho}\) term, this is
the recurrence
\[
 {\cal Q}_{k+1}\le C_p\bigl(1+\lambda^{0.1}\sqrt{{\cal Q}_k}\bigr).
\]
There is a fixed \(K\) with \({\cal Q}_0\lesssim_p\lambda^{-K}\).
After \(n_{\rm buf}\) iterations its potentially growing term is bounded by a
constant times
\(\lambda^{0.2(1-2^{-n_{\rm buf}})-K2^{-n_{\rm buf}}}\).
Choose the fixed \(n_{\rm buf}\)
so that this exponent is positive. This proves bounded ratios on
\(B_{n_{\rm buf}}\), and the same brackets give
\(\||X_i|^2-S_i\|_p\lesssim_p d_0\).

Project both leaves to radius distance comparable to \(d_0\) after
their fork. The synchronized later layers cost
\(\lesssim_{p,\circ}uA^{3/2}/\sqrt L\ll d_0\).
Polarization controls the projected inner product, since the squared
distance of the two branches through their common prefix is
\(O_p(d_0)\). This proves the first assertion of \eqref{eq:C15}.


\paragraph{Second moments in radius and in mass.}

For the second bound, over each box we have
\[
 \int\mathbf EW\,dr\lesssim u^3 A^3d_0,\qquad
 \int\mathbf E\operatorname{tr}C_v^2\,dr\lesssim u^2 A^3d_0.
\]
These estimates follow from \eqref{eq:A1}, the \(S\) oscillation
bound, and constant removal strength. Use earlier and later windows of
length \(h=t'd_0\). The sliding variance estimate proved at
\eqref{eq:C10} is now
\[
 \int_{J_h}\operatorname{Var}(\overline L^{\,\pm}_{r,h})\,dr
 \lesssim\frac1{\kappa_jh}\int_{J^+}\mathbb EW\,dr.
\]
Here \(J_h\) omits endpoint strips of width \(h\), so the windows
stay inside \(J\); \(J^+\) merely provides the available integral
upper bound. Discard those strips with the pointwise moment bound.
Apply the absorption calculation at \eqref{eq:C10} on height scale
\(d_0\), now with trace term \({\cal T}\lesssim A^3/L^2\) and
variance term \((A/L)^3/(\kappa_jt')\). It gives
\[
 \int D^2dr/d_0^3\lesssim t'+(A/L)^3/(\kappa_j t')+A^3/L^2.
\]
Take \(t'=(A/L)^{3/2}/\sqrt{\kappa_j}\le\lambda^{0.18}\).
This proves the radius estimate and also the mass estimate on the floor,
whose density is bounded above. To treat the remaining mass, enlarge
endpoint caps to whole caps. Total mass and \(S\) oscillation remain
\(O(d_0)\). Equation \eqref{eq:C6} and PE (20), with \(p=0\),
give the same integrated \(W\) and trace bounds on caps and contact.

Order this measure by the deterministic coordinate
\(s(v)=\int^v\mathbf1_{\rm nonfloor}(t)\,dt\), and take windows
of \(s\)-length \(h=t'd_0\). The optimizer and its floor set were
fixed before sampling; hence the window indicators are deterministic
weights in \eqref{eq:A1}. Fubini now bounds the integrated centered
variance by \(h^{-1}\int\mathbb EW\,dv\), with no floor-density
denominator. The ordered identity and monotonicity control the noise
and mean changes as in radius. If the whole nonfloor measure has length
less than two windows, the endpoint-strip bound alone suffices. This
proves the second assertion of \eqref{eq:C15} in both measures.


\paragraph{The alias error.}

For the final bound of \eqref{eq:C15}, it remains by \eqref{eq:A1}'s covariance convention to bound \(\mathbf E(X_i\cdot(y-X_i))^2\). Telescope \(X_i\) into past layers ending at dyadic distances from the current radius, down to \(x^{40}\). For an increment ending at distance \(t'\) after starting at distance \(\asymp t'\), with \(x^{40}\lesssim t'\ll u\), \eqref{eq:A3} costs for its squared product with \(y-X_i\) at most \(O_\circ(x^6/(u^2 C_0 t'))\) times the layer's squared norm in expectation. That norm has box average \(O(t'\lambda^{-\rho})\). In radius this follows for \(t'\le d_0\) by sliding the \(S\) variation, for \(d_0\le t'\ll b_{\rm buf}\) pointwise, and otherwise by \eqref{eq:C8}.

In mass, the floor uses the radius estimate, and around a nonfloor evaluation the \(S\) oscillation at distances \(O(t')\) within the buffered patch is \(O(t')\), by \eqref{eq:C6} and the lower potential bound \(-O(C_0/M_j^2)\), as at the good cut. The unprojected last past increment thus has negligible squared norm on average by the same facts, and the initial anchor, projected from fixed comparable distance within constant strength, costs \(O_\circ(u^2 A^3)\). Summing proves the last assertion.

\end{proof}


Here are needed integrated estimates on free forks. Choose a sufficiently large fixed \(c_H\) so \(v\ge c_Hu\) in mass lies beyond the extra floor, and away from the observed buffer with room for prior comparable projection intervals. \begin{lemma}[Integrated errors at free forks]\label{bud:free-forks}
Uniformly,
\[
\begin{split}
 \int_0^{c_Hu}D_v dv&\lesssim_\circ u^2 T_1,\qquad T_1=A+B_f+y_c^2,\\
 \int_0^{c_Hu}D_v^2dv&\lesssim_\circ u^3 T_2,\quad
 T_2=A^2+B_f^2+A^{3/2}y_c+A y_c^{5/3},\\
 (\operatorname{avg}_{v\sim w}D_v^2)^{1/2}&\lesssim_\circ u A
                   \quad (w\ge c_Hu).
\end{split}                                                          \tag{C16}\label{eq:C16}
\]

\end{lemma}\begin{proof}

Regular bands use \eqref{eq:C10}, with \(y_c\gtrsim\sqrt A\); the inner region \(v\lesssim w_c\) uses \eqref{eq:C8} for the first line, also on floor (of mass \(O(u A)\)) for the second.

For contact and comparable-endpoint caps in inner dyads \(v/u\asymp t'\le O(y_c)\), (20)--(21) with constant field slope lower bound \(C_0\) and raw scale \(u y_c\) give for the dyadic mass average of \(D^2/u^2\) at most
\[
 O_\circ((A^3 y_c/t'^2)^{2/3}+A^3/t'^2)
\]
apart from negligible cap-width and strip errors. Clip also by \(O(y_c^2)\); summation weighted by \(t'\) costs \(O_\circ(A^{3/2}y_c)\).

For noncomparable-endpoint caps meeting the region the upper masses are \(O(u y_c)\) by \eqref{eq:C7} and the width bound. Slice each such cap geometrically into units at masses comparable to their lengths, of scale \(u t'\). On one unit the whole-cap \(\Delta S\le M_j^{-1}\) trace proof and constant cap density cost an additional factor \(O(y_c/t')\) in both trace bounds of (20). Use one height bin there: in (21)'s proof the squared mean-oscillation term costs negligible since \(S\) varies by at most \(M_j^{-1}\), including for reverse ordering (its fourth moment estimate uses \eqref{eq:C8}). Thus the same average is bounded by \(O_\circ(A^3 y_c/t'^3)\), again clipped by \(O(y_c^2)\). There are boundedly many such units per mass dyad by disjointness; summation gives \(O_\circ(A y_c^{5/3})\). All band calculations can stop at \(v\sim x^{10}\) with negligible cost. This proves \eqref{eq:C16}.

For completeness, the additional broad-cap cost can be seen directly.
Writing \(y=y_c\), its normalized mass sum is bounded by
\[
 \sum_{t'\,\text{dyadic}}t'\min\{y^2,A^3y/(t')^3\}.
\]
The two terms meet at \(t'_0=A y^{-1/3}\).  Below this scale the
geometric sum is at most \(Cy^2t'_0\); above it the sum is at most
\(CA^3y/(t'_0)^2\).  Both equal \(O(Ay^{5/3})\).
If the crossover is outside the finite summation range, the corresponding
one-sided geometric bound is only smaller.  The original-floor raw
contribution is at most \(O(Ay^2)\), absorbed in this same term since
\(y=o(1)\).  The analogous crossover for comparable caps is
\(t'_0=A^{3/2}/y\), giving \(O(A^{3/2}y)\).

The trace conversion on a sliced cap uses the lower mass bound of that
slice.  The whole-cap numerator is divided by its smaller slice mass,
which is precisely the factor \(y/t'\) already paid above.  No trace
estimate is applied below its lower mass cutoff.  These are mass
estimates; the radius estimates used in Lemma~\ref{bud:box-moments}
and in the coarse profile proof were proved separately by radius slides.
\end{proof}



\begin{lemma}[High projection and centered primitives]\label{bud:high-primitive}
Work in the buffered violation constructed above, with
\(d_0=uL\), \(L\ge a_{\rm goal}Z\),
\(h_1=u\max(L,\sqrt\lambda)\), and
\(G_1=\min(\max(L^2,\lambda),\Lambda)\).
Here \(Z\ge1\) is the deterministic majorant in the global parameter
from Lemma~\ref{bud:envelope}; it is not a random variable under the
positive spin law. For a high dyad \(v\sim w\ge c_Hu\), replace an
overlap at the observed buffer fork by the inner product of the two
branches projected to a fixed cut of order \(w\), after their
separation and strictly before the dyad. In each fixed moment the cost
is at most
\[
 O_{p,\circ}(z_w),\qquad z_w=x^3\min\{(w\sqrt{C(w)})^{-1},\lambda^{-1/2}w^{-2}\}.
                                                                    \tag{C17}\label{eq:C17}
\]
The integral of \(L_v-\mathbb EL_v\) over a high dyad, against any
deterministic weight bounded by one, has \(L^2\) norm
\(o(u^2G_1)\), with a fixed power saving. This includes the plateau
for every edge multiplier used on the comparison path.
\end{lemma}
\begin{proof}
The first option in \eqref{eq:C17} is the postprojection bound obtained
after \eqref{eq:C9}. The actual field gives the second option:
\(dh\gtrsim\lambda w^2\,dv\) on comparable active intervals by
\eqref{eq:P4} and bounded active density; the edge ramp loss is
absorbed there. Even at the plateau or a macroscopic band, a preceding
active interval is available at a comparable scale. Telescope both
branches in regular layers to obtain the stated minimum.

We next estimate the centered primitive. On the active part,
\eqref{eq:A1} gives
\[
 \left\|\int k_v(L_v-\mathbb EL_v)\,dv\right\|_2
 \lesssim_\circ x^3\min\{C(w)^{-1/2},\lambda^{-1/2}/w\}
 \lesssim_\circ u^2X^3\Lambda^{-1/3}\lambda^{-1/6}.
\]
For the first bound integrate \(\mathbb EW\) using PE (20) on
comparable caps and contact, and in radius on the floor. Complete a
partial cap before estimating it. This option also holds on the
plateau. The second option uses the actual-field slope just given;
balancing the two in \(w/u\) gives the last expression. Its ratio to
\(u^2G_1\) is power small by \eqref{eq:C18} below.

For the plateau, first suppose \(\chi_e>0\). From \eqref{eq:C13},
\({\cal N}\le x^2\lambda^{-s_0}Z\le U^2Z\).
On a plateau strip of width \(d'\), Cauchy--Schwarz and monotonicity
compare \(S\) to the constant reference \(q\) on neighboring
intervals of length comparable to \(d'\). The denominator is
\(\sqrt{d'}\), so the resulting oscillation bound is
\[
 \operatorname{osc}S\lesssim {\cal N}/\sqrt{d'}
          \lesssim ZU^2/\sqrt{d'}.
\]
On the capped left strip, use also the preceding active interval of
mass comparable to \(d_-\). Its reference varies by at most
\(O(E_b)=O(U^2/\sqrt{d_-})\), which fits the same bound. On the
capped right strip, boundedness suffices, because its width is
\(d_+=U^4\) and \(U^2/\sqrt{d_+}=1\).

The edge slope on this strip is
\(dh/dv\gtrsim\lambda_1\chi_e U^2/(d')^{3/2}\).
Thus \eqref{eq:A1} bounds the squared primitive norm there by
\[
 O_\circ\left(
  \frac{x^6(d')^{3/2}}{\lambda_1\chi_e U^2}
  \frac{ZU^2}{\sqrt{d'}}\right)
 =O_\circ\left(\frac{x^6Zd'}{\lambda_1\chi_e}\right).
\]
Summing the dyadic strips gives the edge bound
\(O_\circ(x^3\lambda_1^{-1/2}\chi_e^{-1/2}\sqrt Z)\).
When \(\chi_e=1\), the auxiliary-removal bound at \(w\asymp1\)
is \(O_\circ(x^3u^{-3/2}\Lambda^{-1/2})\). Taking the smaller
bound, and using \(\min(a,b)\le a^{2/3}b^{1/3}\), gives
\[
 O_\circ(u^2X^3\Lambda^{-1/3}\lambda_1^{-1/6}Z^{1/6}).
\]
Equation \eqref{eq:C18} again makes its ratio to \(u^2G_1\)
power small, including the dependence on \(Z\).

Other values of \(\chi_e\) occur only at
\(\lambda=x^{p_*}\), with \(\chi_{\rm lin}=1\).
The violation bound supplies the useful lower bound
\[
 u^2G_1\ge u^2\sqrt\lambda L
       \ge x^2\lambda^{1/2-s_0}Z.
\]
There are three cases.
\begin{enumerate}[label=(\alph*)]
\item If \(\chi_e\ge x^{h_e}\), the edge bound divided by this
lower bound is at most
\[
 O_\circ\left(
 x^{1-h_e/2}\lambda_1^{-1/2}\lambda^{s_0-1/2}/\sqrt Z\right)=o(1).
\]
\item If \(\chi_e<x^{h_e}\) but \(u\ge x^{h_e}\), the linear
seed gives primitive norm \(O(x^{3/2})\) by \eqref{eq:A1} and
boundedness of \(S\). Since \(G_1\ge\lambda\), its ratio to the
target is at most
\(O(x^{3/2-2h_e}\lambda^{-1})=o(1)\).
\item If both \(\chi_e<x^{h_e}\) and \(u<x^{h_e}\), remove
strips of width \(\delta=x^{2+h_e/8}\) at the two plateau ends.
Their primitive norm is \(O(\delta)\) by boundedness, and their
ratio to the target is
\[
 O\left(x^{h_e/8}\lambda^{s_0-1/2}/Z\right)=o(1).
\]
On the remaining plateau, the stronger part of \eqref{eq:C13}
and the same neighboring-interval argument give
\[
 \operatorname{osc}S
 \lesssim x^{2+h_e/4}Z/\sqrt\delta
 =x^{1+3h_e/16}Z.
\]
The linear slope is \(dh/dv\gtrsim x^3\), so \eqref{eq:A1}
gives primitive norm at most \(O(x^{2+3h_e/32}\sqrt Z)\).
Its ratio to the target is
\(O(x^{3h_e/32}\lambda^{s_0-1/2}/\sqrt Z)=o(1)\).
\end{enumerate}
All displayed ratios have fixed positive power margins because
\(Z\ge1\), \(h_e\) is small and fixed, and \(p_*\ll h_e\).
This proves the plateau assertion in every case.
\end{proof}


\begin{lemma}[Strict power margins]\label{bud:power-margins}

For clarity, here is some power accounting used here and below (write \(\ll\) in such estimates for polynomial gains). Uniformly under \eqref{eq:C4}, \eqref{eq:C11}, \(a_{\rm goal}Z\le L\le1\), we have
\[
\begin{gathered}
 T_2,\ A T_1,\ X^3 A\Lambda^{-1/3}\lambda^{-1/6},
          \ \lambda^{-2\rho}A^3\ \ll P_1,\qquad
 \lambda^{0.06}T_1,\ L T_1,\ X^3\Lambda^{-1/3}\lambda_1^{-1/6}Z^{1/6}
                 \ \ll G_1 .                                             \tag{C18}\label{eq:C18}
\end{gathered}
\]

\end{lemma}\begin{proof}

Write \(X=\lambda^\xi\), with
\(\xi\ge s_*/2-o(1)\). For terms whose numerator does not depend
on \(L\), monotonicity of \(G_1\) and \(P_1=LG_1\) reduces
the estimate to \(L=a_{\rm goal}\). For \(LT_1/G_1\), also
check \(L=\sqrt\lambda\) and \(L=1\); these are the remaining
possible extrema of the piecewise ratio, and the former requires
\(\xi\ge s_0/2+1/4\). The \(Z\)-dependent term is treated
separately below. Finally \(AT_1\lesssim T_2\); inserting the
powers of \(A\) shows that the \(X^3A\) and
\(\lambda^{-2\rho}A^3\) ratios have larger margins than the
principal ratios computed next, for sufficiently small \(\rho\).

The tables below give the three principal ratio calculations. At
\(s_*=s_0=1/2\), \(\eta=0\), \(k_c=0.03\), and with the
arbitrarily small \(\zeta\)-loss suppressed, write \(e(F)\) for
the power of \(\lambda\) in \(F\). At \(L=a_{\rm goal}\),
\(e(P_1)=2\xi-1/2+e(G_1)\), and the required piecewise functions are
\[
\begin{array}{c|c|c}
 F&\text{range of }\xi&e(F)\\ \hline
 G_1 &[1/4,103/400]&3/100\\
     &[103/400,1/2]&4\xi-1\\
     &[1/2,\infty)&1\\ \hline
 T_1 &[1/4,99/200]&2\xi-1/100\\
     &[99/200,\infty)&3\xi/2+19/80\\ \hline
 T_2 &[1/4,1/2]&11\xi/3-11/600\\
     &[1/2,57/50]&13\xi/4+19/100\\
     &[57/50,\infty)&3\xi+19/40
\end{array}
\]
The first two pieces for \(T_2\) come from \(Ay_c^{5/3}\), and
the last from \(B_f^2\). Subtracting the denominator exponents
shows that all three minima occur at \(\xi=1/2\):
\[
\begin{array}{c|c|c}
 \text{ratio}&\text{term attaining the minimum}&\text{minimum exponent}\\ \hline
 T_2/P_1&Ay_c^{5/3}/P_1&63/200=0.315\\
 \lambda^{0.06}T_1/G_1&\lambda^{0.06}B_f/G_1&19/400=0.0475\\
 X^3\Lambda^{-1/3}\lambda^{-1/6}/G_1
   &X^3\Lambda^{-1/3}\lambda^{-1/6}/G_1&97/300
\end{array}
\]
On the final intervals the relative exponents have positive slopes.
Thus strictly positive margins persist for the prescribed small
\(\tau,\eta,\tau_0\), and for
sufficiently small \(\rho\), \(\zeta/p_*\), and fixed-order moment
interpolation losses, chosen in the stated parameter order.  The number of strengthening stages is already
fixed at this point.  None of these choices requires decreasing \(p_*\)
again.

The dependence on an unbounded \(Z\) in the primitive estimate can
also be checked directly. Put \(Y=X\sqrt Z=\lambda^\xi\) and
\(L_0=Y^2\lambda^{-s_0}\le L\le1\), so \(\xi\ge s_0/2\).
Since \(G_1\) is nondecreasing in \(L\), the same piecewise
calculation gives
\[
 \frac{X^3\Lambda^{-1/3}\lambda_1^{-1/6}Z^{1/6}}{G_1(L)}
 \le Z^{-4/3}\lambda^{d_*},\qquad
 d_* =\frac{97}{300}-\frac32(\tau+\tau_0)-\frac\eta6>\frac14.
\]
The minimum occurs at \(\xi=s_0/2+1/4\). Thus a large majorant
only improves this ratio; no pointwise upper bound on \(Z\) was used.

\end{proof}




\subsection{Quantitative one-row closure}


\begin{proposition}[Buffered one-row closure]\label{bud:closure}

In any of the boxes of \eqref{eq:C15} we claim
\[
 \operatorname{avg}|S_i-\Gamma_s(K_s(i))|=o(u^3 P_1).                     \tag{C19}\label{eq:C19}
\]

\end{proposition}\begin{proof}

Apply \eqref{eq:A7} through degree two at \(K_s\), with the
untilted law, the two observed labels \(a,b\), and test \(1\).
Each mismatch insertion is \((b-C)\delta\). Linear terms vanish
termwise after averaging spins under the actual law.
For the remainders, expand the cavity-spin test to order two, then each
quotient term of degree \(j\le2\) to order \(2-j\), from the
independent-row endpoint to the actual row. Taylor's integral remainder
has three mismatch factors on a fixed finite extension of the tree.
It retains \(\tau_a\tau_b\), either in the interpolated test or in
one deterministic scalar coefficient; every additional cavity spin
occurs in an inserted pair.

After deleting the row, nonnegative spin-only bounds transfer by the
bounded-covariance Gr\"onwall rule, also during the bounded positive
semidefinite row changes below. Deletion itself costs
\(O(m^{-1})\ll u^3P_1\), since
\(P_1\ge\sqrt\lambda\,a_{\rm goal}^2\).
Keep deleted overlaps during sign arguments, so the other mismatch
factors use only the remaining rows. Changing this convention costs
another \(O(m^{-1})\). Moving a revelation cutoff means an affine
positive semidefinite interpolation of the row before applying PE (2).
It does not change the conditioning in a term already being estimated.

The allocation rule after \eqref{eq:A1} supplies all fork measures.
Distinct free vertices away from the observed fork have bounded joint
densities in their mass coordinates, by sequential attachment.
Repeated attachments at one vertex use its single variable.
Ordering restrictions are deterministic and may be retained or dropped
when taking bounds. An inserted pair is \emph{pinned} if its two leaves
split at the observed fork. These give the only split atoms, since
diagonal pairs are omitted. Every other pair is called low or high
according as its split mass is below or above \(c_Hu\).

\begin{enumerate}[label=(\roman*)]

\item Terms with no pinned pair: two low errors cost \(O_\circ(\int_0^{c_Hu}D^2)\) by marginalization and Cauchy--Schwarz, whether at the same vertex or not. One low and one high cost \(O_\circ(u^3 T_1 A)\) using distinct-vertex densities. Any extras here can be dropped.

If all errors are high one gains an extra \(O_p(u)\) low covariance. Indeed turn off the relevant cavity-row revelation below \(c_Hu\), retaining its later cumulative covariance by a jump there. This includes the scalar row law containing \(\tau_a\tau_b\) if that factor occurs in a deterministic coefficient: restoring that law costs \(O(u)\) by PE (2) since \(K_s=O(u)\) there. Products can be treated factor by factor. With low revelation off, the term containing those two spins vanishes by sign flips of a whole branch component (negating the attached row Gaussians if interacting); they are in distinct components and all other pairs are internal.

In an interpolated interacting row this is an admissible PSD interpolation (move the cumulative row-specific field and rest-row couplings, or their affine mixture with the scalar row, to the jump). Spin-only tests on remaining rows are invariant. The new pair for restoring this row costs \(O_p(u)\) at each low split by \eqref{eq:C8} and bounded row transfer, since its coefficient, apart from \(O(1/m)\), is bounded by \(O(|Q_v^{\rm deleted}|+u)\). Take large fixed moments here and slight higher-than-second moments on the two high errors, costing arbitrarily small power losses by boundedness and \eqref{eq:C16}. Old topology and weights remain fixed during restoration. Thus all-high terms cost \(O_\circ(u^3 A^2)\).
\item At degree at least three with pinned errors: for three pinned factors use one \eqref{eq:C15} averaged RMS saving and pointwise higher orders on two, paying \(O(\lambda^{0.07}d_0^3)\) (or with arbitrarily small extra losses); \(\lambda^{0.07} L^2\ll G_1\).

For two pinned and one free factor use the pointwise costs \(d_0^2\) times \(O_\circ(u A+u^2 T_1)\). With one pinned and two free, use a high fixed moment for the pinned error and second moments with slight power loss for the others. The latter two integrated cost at most \(O_\circ(u^2 A^2+u^3 T_2+u^3 A T_1)\). These bounds and the preceding ones are sufficient by \eqref{eq:C18}.
\item At degree two with one pinned error, a low free error costs
\(O_\circ(d_0\lambda^{0.07}u^2T_1)\), using second moments and the
average at the observed fork. For a high free error, fix its dyad
\(v\asymp w\) and project the pinned overlap as in
\eqref{eq:C17}. The error, including the free-vertex mass length, is
\[
 O_\circ(wz_w\,uA),\qquad
 wz_w\lesssim u^2X^3\Lambda^{-1/3}\lambda^{-1/6}.
\]
It is \(o(u^3P_1)\) by \eqref{eq:C18}.

For the projected term we must keep its joint law fixed as \(v\)
varies. Fix the placement topology and all other free vertex
coordinates. Use a projection cut fixed throughout the dyad, prune the
pinned branches at that cut, and retain one path through the high
vertex to its terminal spin. The sharing of these retained data is
independent of \(v\): any sharing through the high fork already
extends beyond the projection cut. All discarded branches have
normalized continuation kernels. Let \(F\) be the projected centered
pinned factor; it is measurable with respect to the retained earlier
data. Conditional on those data and the high prefix, the high pair has
independent continuations with conditional mean \(X_v\). Hence
\[
 \mathbb E[F(Q_v-S_v)]
 =\mathbb E[F(|X_v|^2-S_v)]
 =2\mathbb E[F(L_v-\mathbb EL_v)].
\]
The second equality conditions the retained terminal spin on its prefix.
The scalar coefficients and allocation restrictions are bounded
deterministic weights \(k_v\), and \(\|F\|_2=O(d_0)\).
Lemma~\ref{bud:high-primitive} therefore gives
\[
 \left|\mathbb E\left[F\int k_v\,2(L_v-\mathbb EL_v)\,dv\right]\right|
 \le2\|F\|_2
       \left\|\int k_v(L_v-\mathbb EL_v)\,dv\right\|_2
 =o(d_0u^2G_1)=o(u^3P_1).
\]
This uses unconditional Cauchy--Schwarz on the fixed marginal law.
If other slots attach at the same high vertex, there is still only one
integration in its mass. Fubini and the power margins permit summation
over topologies and dyads.

\end{enumerate}



\paragraph{Two insertions pinned at the observed fork.}

It remains to treat both degree-two pairs pinned. Write each centered overlap as \(m_i+\eta_{\rm pair}\) with \(m_i=|X_i|^2-S_i\) on the common prefix and \(\eta_{\rm pair}=Q_i-|X_i|^2\). Single \(\eta\) terms vanish by conditioning.

For two \(\eta\)'s the coefficients individually gain \(O(u)\): either a third child has entered at the fork (coefficient \(O(u)\)), or both pairs run across the two occupied children. In the latter case the total number of scalar spins, including old spins, in each child is odd; at least one scalar moment factor has a sign cancellation on turning off scalar revelation through this cut (keep later cumulative clock via a jump after the observation, or take its immediate right limit). Indeed the two children then have independent sign symmetry. If using cuts immediately to the right for a fixed placement with the same children, they separate from all further distinct fork vertices eventually; the restoring bound by PE (2) is uniform \(O(u)\). Thus \eqref{eq:C15}'s final bound and \eqref{eq:C18} suffice.

The factor \(\mathbb E m_i^2\) is common to all pinned topologies,
so extract it before summing coefficients. We claim that the remaining
signed scalar coefficient, before the common \((b-C)^2\), is
\(\Gamma_s''(K_s(i))/2\) at almost every buffered cut.

To identify it, compress a small neighborhood of the observed mass
into a clock jump at that mass, retaining the observation at its
original intermediate clock value. At a continuity cut there is room
to move this value in both directions: \(K_s\) is strictly increasing
locally because its field contains \(C_0r\). On an already extended
old tree, increase the common subincrement at the distinguished vertex
by variance \(\epsilon\), and decrease the following subincrement
in every occupied child by the same amount. The covariance direction
is then one for pairs entering different children of that vertex and
zero for same-child pairs, diagonals, and pairs outside its subtree.
Total leaf variances and all other split covariances stay fixed.
Thus this is exactly the pinned direction, even when earlier
insertions have occupied more than two children.

In the quotient expansion of \eqref{eq:A5}, the degree-two scalar
terms are
\[
 N_2-N_1\Psi_1+N_0(\Psi_1^2-\Psi_2).
\]
In every term with positive denominator degree, allocate one whole
such \(\Psi\) moment last. Previously placed factors pull back from
the induced old genealogy, and the extracted prefix expectation is
independent of placement. The complete signed sum of that denominator
moment is therefore a covariance derivative of the constant-one
observation on this genealogy, hence zero. This cancellation requires
the full signed sum; it is not a row-sum identity for a nonconstant
restricted test. Coincidences within the permitted children remain
part of that sum.

Only \(N_2\) remains. Its coefficient includes the Taylor factor
\(1/2!\). The two subincrements at the same positive mass compose
into the original scalar transition, so moving their intervening split
changes only the evaluation time of \(\Gamma_s\). The coefficient
is consequently \(\Gamma_s''/2\). This calculation introduces no
additional split atom: a new attachment at a same-mass subdivision
without old divergence has zero coefficient, because that edge has
zero length and one continuing child. A different vertex with the
same numerical mass likewise creates no atom at the observed fork.

Finally shrink the compressed neighborhood. Bounded-clock path limits
with the cut retained pass the aggregated coefficients to the original
path. Other freely varying fork masses in the shrinking interval have
vanishing absolute weight. The scalar identity
\(\Gamma_s''=\mathbb EV_{zzz}^2-2\alpha_t\mathbb EV_{zz}^3\)
uses exactly the fork mass at the intermediate observation. At the
limit this remains its mass, since \(K_s\) grows strictly locally
and the spatial derivatives and path moments converge. Null contact
and jump exceptions do not affect either box integral. This proves
the claimed coefficient identity.

By the drift estimates at \eqref{eq:C14}, \(\Gamma_s''(K_s(i))=O(d_0+\lambda u)\). The \(m_i^2\) average by \eqref{eq:C15} is at most \(O(d_0^2\lambda^{0.14})\); their product suffices since \(\lambda^{0.14}L(L+\lambda)\ll G_1\). This completes \eqref{eq:C19}.

In particular the error budget remains relative to the removal strength:
\[
 \frac{u^3P_1}{d_0h_1^2}
   =\frac{G_1}{\max(L^2,\lambda)}
   \asymp\min(1,R_n)
\]
with the notation used next.  This observation will be needed when
\(R_n\) tends to zero; an absolute \(o(1)\) error would not suffice.
\end{proof}




\subsection{Buffered rigidity and completion of the budgets}


\begin{proof}[Proof of Lemma~\ref{bud:strengthening}]
The ratio \(R_n\) defined below compares the removal strength with the
curvature scale of the rescaled box. We first treat \(R_n\to\infty\),
with a second rescaling if the box shrinks, and then bounded \(R_n\),
including \(R_n\to0\). The closure error is
\(o(\min(1,R_n))\); its relative form is essential when calibration
at support points requires division by \(R_n\).


Use \(r=r_0+h_1 z,\ \theta=d_0/h_1\) (local notation, unrelated here to size-sweep time), and indices \(n\) along the growing systems. Put
\[
\begin{aligned}
 m_n(z)&=(\alpha-\beta_0)(r)/d_0,\qquad
 \widetilde m_n(z)=(\widetilde\alpha-\beta_0)(r)/d_0,\\
 Y_n(z)&=(S-r)/d_0,\qquad
 \varphi_n(z)=(\widetilde r(\alpha(r))-r_0)/h_1,\\
 g_n(z)&=[\Gamma_s(k_*(r))-r]/(d_0 h_1^2).
\end{aligned}
\]
The first three functions are bounded on compact sets with a common bound, and we can take \(\beta_0'\to c_\beta>0,\ c_q=1/c_\beta\) in these local ranges. We have on compacts
\[
 g_n''=-A_*\widetilde m_n+o(1),\qquad \varphi_n=z+\theta(Y_n+o(1)),
\]
the curvature statement a.e. (also distributionally, since the scalar first derivatives of \(\Gamma\) are absolutely continuous) with \(A_*\) now its positive limiting value. Here \(\varphi_n\) is a nondecreasing map. Indeed Taylor expand \(\widetilde r(v)\) at \(K_*(v),\ v=\alpha(r)\); the argument increment is \(b(S-q)-C_0(S-r)+d\) and \(\Gamma_*''=O(\lambda u)\), so the quadratic error is \(O(\lambda u d_0^2)=o(u^3P_1)\). Also \(q(v)-r=d_0(c_qm_n+o(1))\). Hence \eqref{eq:C19}, with \(d=o(u^3P_1)\), gives
\[
\begin{aligned}
 g_n(\varphi_n)&=R_n Y_n + D_n(Y_n-c_qm_n)+e_n,\\
 R_n&=\Gamma_*'(K_*(v)) C_0/h_1^2,\qquad
 D_n=[1-b(v)\Gamma_*'(K_*(v))]/h_1^2 .
\end{aligned}
                                                                    \tag{C20}\label{eq:C20}
\]
where \(e_n=o(\min(1,R_n))\) in box-averaged absolute value (both measures as before) on all moving boxes of length comparable to \(\theta\) in compact \(z\)-ranges. By \eqref{eq:P2}--\eqref{eq:P5},
\(D_n\asymp\lambda u^2/h_1^2\lesssim1,\ D_n=o(R_n)\); the pointwise lift pays the small \(v^2\) offset in \eqref{eq:P3}, and \(U\lesssim u\). On compacts \(D_n\) varies by \(o(1)\), \(R_n\) by \(o(R_n)\), by the active derivative bounds (\(b'=O(\lambda u),\ \Gamma_*''=O(\lambda u)\)). We subselect repeatedly.

We have uniform limiting bounds on \(g_n,g_n'\) independent of compact interval. Indeed intervals of length a sufficiently large fixed multiple of \(\theta\) contain nonfloor physical mass \(\gtrsim d_0\) by the buffer bounds on \(m_n\); there \(R_n|Y_n|\le O(R_n/(d_0 M_j))=o(1)\). Thus \eqref{eq:C20} calibrates bounded values at bounded spacings, sufficient by the curvature bound. Take a \(C^1_{\rm loc}\) limit \(g\). Recall \({\cal H}'=2C_0(r-S)\) in radius and
\[
 {\cal H}(\alpha'-\beta_0')\le0,\qquad {\cal H}\ge-O(C_0/M_j^2)
\]
locally since \(\beta_0'\) lies strictly between floor and cap. On nonfloor \({\cal H}\le0\).

\paragraph{Unbounded rescaled removal.}


Suppose first \(R_n\to\infty\). Equation \eqref{eq:C20} gives \(Y_n=o(1)\) in box-averaged radius, hence uniformly on compacts (a violation persists over a fraction of a box on one side, by monotonicity of \(S\)). Then \(\widetilde m_n-m_n\to0\) in local \(L^1\): \(\widetilde r(\alpha(r))-r=o(d_0)\), and by Fubini the area between distributions on each such interval of length \(O(h_1)\) costs \(o(d_0)\) times \(O(h_1)\) mass (crossings are local by monotonicity). Write
\[
 J_n(z)=\Gamma_*'(K_*(\alpha(r_0)))\,{\cal H}(r)/(2d_0 h_1^3),\qquad
 \ell_1=\lim c_q D_n .
\]
The variational sign gives \(J_nm_n'\le0\). Equation
\eqref{eq:C20} gives
\(J_n'=-R_n(0)Y_n=-g_n-\ell_1m_n+o(1)\), where the error tends
to zero in the moving radius-box integrals. Nearby nonfloor
calibrations give \(J_n=O(\theta)\), with a common local bound;
the normalized \(M_j^{-2}\) potential error is \(o(\theta)\).
Take a locally uniform limit \(J\) and a weak-star limit
\(m_n\rightharpoonup m_\infty\). Then
\(J'=-g-\ell_1m_\infty\) and \(g''=-A_*m_\infty\).

For clarity, the energy estimate follows by two integrations by parts.
On a fixed interval, integrate \(J_nm_n'\le0\), substitute the
equation for \(J_n'\), and use
\(m_n=-g_n''/A_*+o(1)\). This gives
\[
 0\ge[J_nm_n]+\int g_nm_n+\ell_1\int m_n^2+o(1),
 \qquad
 \int g_nm_n
 =-[g_ng_n'/A_*]+\int g_n'^2/A_*+o(1).
\]
Consequently
\[
 \int (g_n'^2/A_*+\ell_1m_n^2)\le [g_n g_n'/A_* -J_n m_n]_{\rm left}^{\rm right}+o(1).
\]
All boundary factors are uniformly bounded, independently of the fixed
interval. Passing to the limit and then enlarging the interval first
shows that \(g'\in L^2(\mathbb R)\). Bounded curvature makes
\(g'\) uniformly continuous, hence it tends to zero at both ends.
If \(\theta\to0\), then \(J=0\), so both boundary terms vanish
as the endpoints tend to infinity. If \(\theta\) has positive
limit, the measure
\(d(c_\beta z+(\lim\theta)m_\infty)\) has support unbounded in
both directions. Nonfloor approximants give \(J=0\) on that support,
because the floor density vanishes. Send the endpoints to infinity
through its support; again both boundary terms vanish. The energy
estimate now gives \(g'=0\), and \(g''=-A_*m_\infty\) gives
\(m_\infty=0\). Finally \(J'=-g\) and boundedness of \(J\)
give \(g=0\). For positive limiting \(\theta\), monotonicity of
\(\alpha\) then gives the pointwise contradiction
\(|m_n(0)|=1\).


\paragraph{Unbounded removal: the shrinking-box subcase.}

If \(\theta\to0\), we have \(\ell_1>0\) since \(h_1^2=\lambda u^2\). Put \(K_n(t)=J_n(\theta t)/\theta,\ M_n(t)=m_n(\theta t)\). They are bounded on compacts with uniform limiting bounds, \(K_n'=-\ell_1 M_n+o_{L^1_{\rm loc}}(1)\) by the moving box estimate and \(g=0\), \(K_n M_n'\le0\). \(M_n\) converge on subselection in local \(L^1\) and at continuity points by monotonicity with the deterministic trend added. Now \(\ell_1\int M_n^2\le-[K_n M_n]+o(1)\) on fixed intervals. The limit thus has finite global squared norm; sending endpoints to infinities through continuity points where it tends to zero makes that norm zero. Again monotonicity gives the pointwise contradiction.


\paragraph{Bounded rescaled removal, including a vanishing limit.}

Suppose now that \(R_n\) is bounded. Then \(\theta=1\).
After adding their deterministic trends, \(m_n,Y_n\) are monotone;
subselection therefore gives almost-everywhere and local \(L^1\)
limits. The transport estimate also gives
\(\widetilde m_n-m_n\to0\) in local \(L^1\). Indeed on nonfloor
points the displacement \(\widetilde r(\alpha(r))-r\) is
\(o(d_0)\); floor crossings travel only \(O(d_0)\) and have
density \(o(1)\). Fubini bounds the area between the inverse
distributions on a compact interval by \(o(d_0h_1)\).
It follows that \(g''=-A_*m_\infty\).
The normalized potential
\(\mathcal H(r)/(2C_0d_0h_1)\) has derivative \(-Y_n\).
Its lower bound and nearby nonfloor references give a bounded
nonnegative limiting function \(\bar J\).

We next calibrate at the support of
\(\mu=d(c_\beta z+m_\infty)\). The measures
\(\mu_n=d\alpha/h_1\) in this coordinate converge vaguely to
\(\mu\). Every fixed neighborhood of a support point has positive
limiting mass. Its floor mass is \(o(1)\), while the mass integral
of \(|e_n|/R_n\) tends to zero. Hence it contains a nonfloor point
with \(e_n/R_n\to0\). At this point \(Y_n=o(1)\) and
\(D_n/R_n\to0\), so \eqref{eq:C20} gives
\(g_n(\varphi_n)=o(R_n)\). Shrink the fixed neighborhoods and
choose a diagonal subsequence. This produces approximants to each
support endpoint used below. At a contact the potential is zero; in a
cap its nearest zero is within \(M_j^{-1}\), whose normalized error
vanishes. We conclude that \(\bar J=g=0\) on
\(\operatorname{supp}\mu\), even if \(R_n\to0\).

A floor interval can still have positive rescaled radius length; only
its mass has vanished. We therefore retain the separate radius estimate
for \(e_n/R_n\) on any such support gap. After a further subsequence,
for almost every radius coordinate \(z\),
\[
 s(z)=\lim\varphi_n(z)=z+Y(z),\qquad
 g(s(z))=(\lim R_n)Y(z),\qquad \bar J'=-Y.
\]
The function \(s\) is nondecreasing. On a support gap \((a',b')\),
monotonicity and the endpoint approximants give \(s\in[a',b']\).
The gap is bounded, and \(g''\) is affine there with slope
\(A_*c_\beta>0\).

Suppose \(g'(a')>0\). If \(\lim R_n>0\), then
\(g(s)=(\lim R_n)(s-z)\) forces \(s(z)>z\) almost everywhere
just to the right of \(a'\): this holds if \(s\) is near
\(a'\), where \(g\) is positive to its right, and also if \(s\)
is farther inside the gap. Thus \(\bar J'=-Y<0\), contradicting
\(\bar J(a')=0\) and \(\bar J\ge0\).

If \(R_n\to0\), then \(g(s)=0\), and the sign requires the
relative error. Fix an almost-everywhere point \(z>a'\) where the
radius error divided by \(R_n\) tends to zero. If \(s(z)=a'\),
take nonfloor approximants \(z_n^a\to a'\) as above. Eventually
\(z_n^a<z\), so \(\varphi_n(z_n^a)\le\varphi_n(z)\).
Both images lie near \(a'\), where \(g_n'>0\). Since the value at
the approximant is \(o(R_n)\) and \(R_n\) is asymptotically
constant on compact intervals,
\[
 \liminf_n\frac{g_n(\varphi_n(z))}{R_n(z)}\ge0.
\]
Divide \eqref{eq:C20} by \(R_n\). It gives \(Y(z)\ge0\),
contrary to \(Y(z)=a'-z<0\). Every other zero of \(g\) is
bounded away from \(a'\) in this neighborhood, and hence gives
\(s(z)>z\). Again \(Y>0\) almost everywhere just to the right,
contradicting the potential sign. Reflection excludes
\(g'(b')>0\) in both cases.

But \(g(a')=g(b')=0\) and the affine curvature imply
\[
 g'(a')+g'(b')=A_*c_\beta(b'-a')^2/6>0.
\]
At least one endpoint derivative must be positive. Thus no support gap
exists.

Finally, the boundedness of \(m_\infty\) makes the support of
\(d(c_\beta z+m_\infty)\) unbounded in both directions.  With no gaps
it is all of \(\mathbb R\), so \(g=0\) and \(m_\infty=0\).
In each regime where the normalized discrepancies converge locally in
\(L^1\) to zero, the contradiction at the originally selected point is
pointwise: monotonicity of the inverse distribution, after adding back
its deterministic trend, gives for fixed continuity points \(\pm h\)
\[
 -C h+o(1)\le m_n(0)\le C h+o(1).
\]
Here one may take \(C=c_\beta/\theta_\infty\) when
\(\theta\to\theta_\infty>0\), and \(C=c_\beta\) in the bounded
removal case \(\theta=1\).
Letting \(h\downarrow0\) contradicts \(|m_n(0)|=1\).
For the second-rescaling sequence the identical argument is applied to
\(M_n\), whose deterministic trend has slope tending to \(c_\beta\).
No pointwise relative error at the selected point is used.


This proves \eqref{eq:C12}, completing the finite induction. \end{proof}

\begin{proof}[Completion of Proposition~\ref{bud:targets}]
At its last stage the radius assertion converts by \eqref{eq:C7} to \eqref{eq:C3}; the averaged excess \(O(x^{4-2\zeta}u\Lambda^{1/3})\) suffices there. This completes the local comparison required for \eqref{eq:C1}.

Here are the final exponent checks.  The local pressure derivative has
scale \(\lambda u\) when it is used to bound
\(u\int_{v\asymp u}A_v\,dv\).  The optimized excess therefore contributes
\[
 x^{4-2\zeta}\lambda^{k_c/3-1}
 =x^4\lambda^{-2s_0}
       \bigl(x^{-2\zeta}\lambda^{k_c/3-1+2s_0}\bigr).
\]
The parenthesized factor is power small because
\(k_c/3>1-2s_0\) and \(\zeta\) is chosen after \(p_*\).
The discrepancy term in \eqref{eq:C3} contributes
\(u^4a_{\rm goal}^2=x^4\lambda^{-2s_0}\).
Since \(U^4=x^4\lambda^{-2s_*}\), the remaining factor
\(\lambda^{2\tau_0}\) pays dyadic logarithms and the permitted small
losses.  At the smallest dyad \(u\asymp U\), the switch error from
\eqref{eq:P11}, divided by \(\lambda u\) and then by \(U^4\), is at most
\[
 x^{c'}\lambda^{5s_*/2-1},
\]
which has a strict saving.  Larger dyads only improve this estimate.
The global budget handles the remaining dyads
\(u\ge\lambda^{b_\ell}\), as in the initial reduction.
Both parts of \eqref{eq:C2} were already obtained from the global
comparison, with a final-strength power independent of subsequent
sufficiently small choices of \(p_*\).  Thus the local induction does
not enter the choice of the later size Taylor order.
\end{proof}
